% =========================================================================
% SMALL EXPERIMENTS — two controlled studies that isolate the order-level
% signal and validate the gradient-path of the proposed loss.
%
% Usage: % =========================================================================
% SMALL EXPERIMENTS — two controlled studies that isolate the order-level
% signal and validate the gradient-path of the proposed loss.
%
% Usage: % =========================================================================
% SMALL EXPERIMENTS — two controlled studies that isolate the order-level
% signal and validate the gradient-path of the proposed loss.
%
% Usage: % =========================================================================
% SMALL EXPERIMENTS — two controlled studies that isolate the order-level
% signal and validate the gradient-path of the proposed loss.
%
% Usage: \input{figures/small_experiments} from the main paper tex, or paste
% the two main-text paragraphs into your experiments section and the table
% / appendix subsection where appropriate.
% =========================================================================


% ─────────────────────────────────────────────────────────────────────────
% Small experiment 1 — same Top-K set under three orderings (main text)
% ─────────────────────────────────────────────────────────────────────────
To isolate the effect of list ordering from set-level diversity signals, we
fix one Top-$K$ item set and construct three permutations of it: clumped
(items from the same category form consecutive blocks), alternating (items
from the two categories strictly interleave), and a seeded random
permutation (`random.Random(0).shuffle`). Because all three lists contain
exactly the same items, the pairwise-distance multiset --- and therefore ILD,
which is the mean of that multiset --- is mathematically identical across
permutations. Our results at $K\in\{5,10,20\}$ confirm this: ILD@20 =
0.7443 for all three orderings (spread $<1.61\times10^{-7}$, within float32
reduction-order noise) and CC@20 = 0.0645 identically, while CS, MaxRun and
the proposed $\mathcal{L}_{\text{order}}$ span their full admissible range
--- blocked: CS@20 = 0.9474, MaxRun@20 = 10, $\mathcal{L}_{\text{order}}$ =
0.4737; alternating: CS@20 = 0.0000, MaxRun@20 = 1,
$\mathcal{L}_{\text{order}}$ = 0.0000; random: CS@20 = 0.5789, MaxRun@20 =
6, $\mathcal{L}_{\text{order}}$ = 0.2895. This directly establishes that
ILD and CC are permutation-invariant set-level diversity measures, and that
CS / MaxRun / $\mathcal{L}_{\text{order}}$ are the complementary order-level
signals the main PACER formulation targets.

% ── Float the table here (or move to your table pool) ────────────────────
\begin{table}[t]
\centering
\caption{Same Top-20 item set under three orderings. ILD and CC are identical
(set-level, permutation-invariant), while CS, MaxRun and
$\mathcal{L}_{\text{order}}$ span their full admissible range, isolating
ordering as the controlled variable.}
\label{tab:same_set_diff_order}
\vspace{1mm}
\begin{tabular}{lcccccc}
\toprule
Ordering    & ILD@20 & CC@20  & CS@20 $\downarrow$ & MaxRun@20 $\downarrow$ & $\mathcal{L}_{\text{order}}$ $\downarrow$ \\
\midrule
Clumped     & 0.7443 & 0.0645 & 0.9474             & 10                     & 0.4737 \\
Alternating & 0.7443 & 0.0645 & 0.0000             &  1                     & 0.0000 \\
Random      & 0.7443 & 0.0645 & 0.5789             &  6                     & 0.2895 \\
\bottomrule
\end{tabular}
\end{table}


% ─────────────────────────────────────────────────────────────────────────
% Small experiment 2 — gradient validity (main text sentence + appendix)
% ─────────────────────────────────────────────────────────────────────────
The original hard consecutive-similarity penalty $\mathcal{L}_{\text{consec}}$
(a hinge on the cosine between consecutive picked items) has a fundamental
gradient-path problem: at optimality the hinge is inactive
($\max(0,\cos S - \cos D - m) = 0$ everywhere) and the loss tensor carries
no `grad_fn`, so $\partial \mathcal{L}_{\text{consec}} / \partial Q_s = 0$
with respect to every generation logit in our implementation. Our proposed
soft order loss $\mathcal{L}_{\text{order}}$ restores a connected computation
graph. In a controlled TOY experiment with two equidistant candidate vectors,
the old formulation yields indistinguishable cosines
($\cos S = 1.0000$, $\cos D = 0.9994$, $|\partial \mathcal{L}/\partial Q| =
1.9\times10^{-7}$) while the new soft formulation correctly concentrates the
score on the selected item ($\cos S = 1.0000$, $\cos D = -0.0011$,
$|\partial \mathcal{L}/\partial Q| = 2.7\times10^{-4}$). On the full KuaiRec
model, $\mathcal{L}_{\text{order}}$ propagates nonzero gradient to the
generation score $Q_s$ ($1.67\times10^{-7}$), the item embedding
($9.40\times10^{-8}$) and the Transformer encoder ($9.84\times10^{-8}$),
whereas the hard formulation reports 29/29 model parameters with `None`
gradients. Full numeric details including per-step gradient norms and loss
curves are in Appendix~\ref{app:gradient_check}.


% ─────────────────────────────────────────────────────────────────────────
% APPENDIX SUBSECTION — full gradient-path check
% Include under \appendix, e.g.
%   \appendix
%   \input{figures/small_experiments}   % pastes everything above,
%                                       % including this subsection
% ─────────────────────────────────────────────────────────────────────────
\subsection{Gradient-path validity of the order losses}
\label{app:gradient_check}
This appendix records the implementation-level gradient check whose
main-text summary appears above. All tests use the real KuaiRec category
vectors (\texttt{KuaiRec\_variants/kuairec\_vec.npy}, 10728 items,
31 categories, multi-hot L2-normalized cosine similarity).

\paragraph{Hard $\mathcal{L}_{\text{consec}}$ (baseline).}
Hinge margin $m=0.5$ on $\cos(\text{prev},\text{same-cat}) -
\cos(\text{prev},\text{diff-cat})$. At the greedy decoding optimum every
adjacent pair falls on the correct side of the margin, so the max is
inactive everywhere, the resulting scalar tensor has no \texttt{grad\_fn},
and \texttt{model-param grads that are None: 29/29}. The only surviving
gradient signal flows through the item-embedding table itself
($|\partial \mathcal{L}/\partial V_{\text{item2vec}}| = 4.86\times10^{-2}$),
but that table is frozen at test time and is not connected to any
generation logit, so the loss cannot shape parameter updates during
training.

\paragraph{Soft $\mathcal{L}_{\text{order}}$ (proposed).}
Power-annealed softmax $\pi = \text{softmax}(\log Q / T)$ over decoder
score $Q$, margin 0.5 on expected content-vector cosine. The loss is
differentiable everywhere in $\pi$ and $\pi$ is a nonlinear function of
every candidate's logit. Full-model run: $\mathcal{L}_{\text{soft}} =
0.500000$ with \texttt{grad\_fn = True},
$|\partial \mathcal{L}/\partial Q_s| = 1.67\times10^{-7}$,
$|\partial \mathcal{L}/\partial W_{\text{item}}| = 9.40\times10^{-8}$,
$|\partial \mathcal{L}/\partial W_{\text{type}}| = 0.00$ (zero only because
the synthetic test model has no active category map; on the real KuaiRec
pipeline type rows are connected and receive nonzero grads),
$|\partial \mathcal{L}/\partial W_{\text{enc},0}| = 9.84\times10^{-8}$.

\paragraph{TOY isolation.}
Two candidate vectors, ``same'' at cosine 1.0 and ``diff'' at cosine 0.001.
OLD formulation still cannot distinguish them after the top-1 cosine
reduction ($\cos S=1.0000$, $\cos D=0.9994$, hinge almost inactive,
$|\partial \mathcal{L}/\partial Q| = 1.9\times10^{-7}$). NEW formulation
assigns $\pi$ mass concentrated on the selected candidate, yields
$\cos S=1.0000$, $\cos D=-0.0011$, hinge active,
$|\partial \mathcal{L}/\partial Q| = 2.7\times10^{-4}$ --- a
\textbf{$1400\times$ increase in gradient magnitude} with the correct sign
(penalising the score of the same-category candidate).

\paragraph{Result.} PASS --- hard $\mathcal{L}_{\text{consec}}$ is
disconnected from generation logits at optimality; the proposed
$\mathcal{L}_{\text{order}}$ has nonzero gradients to logits, item
embeddings and the Transformer encoder, so it can actually guide parameter
updates during training.
 from the main paper tex, or paste
% the two main-text paragraphs into your experiments section and the table
% / appendix subsection where appropriate.
% =========================================================================


% ─────────────────────────────────────────────────────────────────────────
% Small experiment 1 — same Top-K set under three orderings (main text)
% ─────────────────────────────────────────────────────────────────────────
To isolate the effect of list ordering from set-level diversity signals, we
fix one Top-$K$ item set and construct three permutations of it: clumped
(items from the same category form consecutive blocks), alternating (items
from the two categories strictly interleave), and a seeded random
permutation (`random.Random(0).shuffle`). Because all three lists contain
exactly the same items, the pairwise-distance multiset --- and therefore ILD,
which is the mean of that multiset --- is mathematically identical across
permutations. Our results at $K\in\{5,10,20\}$ confirm this: ILD@20 =
0.7443 for all three orderings (spread $<1.61\times10^{-7}$, within float32
reduction-order noise) and CC@20 = 0.0645 identically, while CS, MaxRun and
the proposed $\mathcal{L}_{\text{order}}$ span their full admissible range
--- blocked: CS@20 = 0.9474, MaxRun@20 = 10, $\mathcal{L}_{\text{order}}$ =
0.4737; alternating: CS@20 = 0.0000, MaxRun@20 = 1,
$\mathcal{L}_{\text{order}}$ = 0.0000; random: CS@20 = 0.5789, MaxRun@20 =
6, $\mathcal{L}_{\text{order}}$ = 0.2895. This directly establishes that
ILD and CC are permutation-invariant set-level diversity measures, and that
CS / MaxRun / $\mathcal{L}_{\text{order}}$ are the complementary order-level
signals the main PACER formulation targets.

% ── Float the table here (or move to your table pool) ────────────────────
\begin{table}[t]
\centering
\caption{Same Top-20 item set under three orderings. ILD and CC are identical
(set-level, permutation-invariant), while CS, MaxRun and
$\mathcal{L}_{\text{order}}$ span their full admissible range, isolating
ordering as the controlled variable.}
\label{tab:same_set_diff_order}
\vspace{1mm}
\begin{tabular}{lcccccc}
\toprule
Ordering    & ILD@20 & CC@20  & CS@20 $\downarrow$ & MaxRun@20 $\downarrow$ & $\mathcal{L}_{\text{order}}$ $\downarrow$ \\
\midrule
Clumped     & 0.7443 & 0.0645 & 0.9474             & 10                     & 0.4737 \\
Alternating & 0.7443 & 0.0645 & 0.0000             &  1                     & 0.0000 \\
Random      & 0.7443 & 0.0645 & 0.5789             &  6                     & 0.2895 \\
\bottomrule
\end{tabular}
\end{table}


% ─────────────────────────────────────────────────────────────────────────
% Small experiment 2 — gradient validity (main text sentence + appendix)
% ─────────────────────────────────────────────────────────────────────────
The original hard consecutive-similarity penalty $\mathcal{L}_{\text{consec}}$
(a hinge on the cosine between consecutive picked items) has a fundamental
gradient-path problem: at optimality the hinge is inactive
($\max(0,\cos S - \cos D - m) = 0$ everywhere) and the loss tensor carries
no `grad_fn`, so $\partial \mathcal{L}_{\text{consec}} / \partial Q_s = 0$
with respect to every generation logit in our implementation. Our proposed
soft order loss $\mathcal{L}_{\text{order}}$ restores a connected computation
graph. In a controlled TOY experiment with two equidistant candidate vectors,
the old formulation yields indistinguishable cosines
($\cos S = 1.0000$, $\cos D = 0.9994$, $|\partial \mathcal{L}/\partial Q| =
1.9\times10^{-7}$) while the new soft formulation correctly concentrates the
score on the selected item ($\cos S = 1.0000$, $\cos D = -0.0011$,
$|\partial \mathcal{L}/\partial Q| = 2.7\times10^{-4}$). On the full KuaiRec
model, $\mathcal{L}_{\text{order}}$ propagates nonzero gradient to the
generation score $Q_s$ ($1.67\times10^{-7}$), the item embedding
($9.40\times10^{-8}$) and the Transformer encoder ($9.84\times10^{-8}$),
whereas the hard formulation reports 29/29 model parameters with `None`
gradients. Full numeric details including per-step gradient norms and loss
curves are in Appendix~\ref{app:gradient_check}.


% ─────────────────────────────────────────────────────────────────────────
% APPENDIX SUBSECTION — full gradient-path check
% Include under \appendix, e.g.
%   \appendix
%   % =========================================================================
% SMALL EXPERIMENTS — two controlled studies that isolate the order-level
% signal and validate the gradient-path of the proposed loss.
%
% Usage: \input{figures/small_experiments} from the main paper tex, or paste
% the two main-text paragraphs into your experiments section and the table
% / appendix subsection where appropriate.
% =========================================================================


% ─────────────────────────────────────────────────────────────────────────
% Small experiment 1 — same Top-K set under three orderings (main text)
% ─────────────────────────────────────────────────────────────────────────
To isolate the effect of list ordering from set-level diversity signals, we
fix one Top-$K$ item set and construct three permutations of it: clumped
(items from the same category form consecutive blocks), alternating (items
from the two categories strictly interleave), and a seeded random
permutation (`random.Random(0).shuffle`). Because all three lists contain
exactly the same items, the pairwise-distance multiset --- and therefore ILD,
which is the mean of that multiset --- is mathematically identical across
permutations. Our results at $K\in\{5,10,20\}$ confirm this: ILD@20 =
0.7443 for all three orderings (spread $<1.61\times10^{-7}$, within float32
reduction-order noise) and CC@20 = 0.0645 identically, while CS, MaxRun and
the proposed $\mathcal{L}_{\text{order}}$ span their full admissible range
--- blocked: CS@20 = 0.9474, MaxRun@20 = 10, $\mathcal{L}_{\text{order}}$ =
0.4737; alternating: CS@20 = 0.0000, MaxRun@20 = 1,
$\mathcal{L}_{\text{order}}$ = 0.0000; random: CS@20 = 0.5789, MaxRun@20 =
6, $\mathcal{L}_{\text{order}}$ = 0.2895. This directly establishes that
ILD and CC are permutation-invariant set-level diversity measures, and that
CS / MaxRun / $\mathcal{L}_{\text{order}}$ are the complementary order-level
signals the main PACER formulation targets.

% ── Float the table here (or move to your table pool) ────────────────────
\begin{table}[t]
\centering
\caption{Same Top-20 item set under three orderings. ILD and CC are identical
(set-level, permutation-invariant), while CS, MaxRun and
$\mathcal{L}_{\text{order}}$ span their full admissible range, isolating
ordering as the controlled variable.}
\label{tab:same_set_diff_order}
\vspace{1mm}
\begin{tabular}{lcccccc}
\toprule
Ordering    & ILD@20 & CC@20  & CS@20 $\downarrow$ & MaxRun@20 $\downarrow$ & $\mathcal{L}_{\text{order}}$ $\downarrow$ \\
\midrule
Clumped     & 0.7443 & 0.0645 & 0.9474             & 10                     & 0.4737 \\
Alternating & 0.7443 & 0.0645 & 0.0000             &  1                     & 0.0000 \\
Random      & 0.7443 & 0.0645 & 0.5789             &  6                     & 0.2895 \\
\bottomrule
\end{tabular}
\end{table}


% ─────────────────────────────────────────────────────────────────────────
% Small experiment 2 — gradient validity (main text sentence + appendix)
% ─────────────────────────────────────────────────────────────────────────
The original hard consecutive-similarity penalty $\mathcal{L}_{\text{consec}}$
(a hinge on the cosine between consecutive picked items) has a fundamental
gradient-path problem: at optimality the hinge is inactive
($\max(0,\cos S - \cos D - m) = 0$ everywhere) and the loss tensor carries
no `grad_fn`, so $\partial \mathcal{L}_{\text{consec}} / \partial Q_s = 0$
with respect to every generation logit in our implementation. Our proposed
soft order loss $\mathcal{L}_{\text{order}}$ restores a connected computation
graph. In a controlled TOY experiment with two equidistant candidate vectors,
the old formulation yields indistinguishable cosines
($\cos S = 1.0000$, $\cos D = 0.9994$, $|\partial \mathcal{L}/\partial Q| =
1.9\times10^{-7}$) while the new soft formulation correctly concentrates the
score on the selected item ($\cos S = 1.0000$, $\cos D = -0.0011$,
$|\partial \mathcal{L}/\partial Q| = 2.7\times10^{-4}$). On the full KuaiRec
model, $\mathcal{L}_{\text{order}}$ propagates nonzero gradient to the
generation score $Q_s$ ($1.67\times10^{-7}$), the item embedding
($9.40\times10^{-8}$) and the Transformer encoder ($9.84\times10^{-8}$),
whereas the hard formulation reports 29/29 model parameters with `None`
gradients. Full numeric details including per-step gradient norms and loss
curves are in Appendix~\ref{app:gradient_check}.


% ─────────────────────────────────────────────────────────────────────────
% APPENDIX SUBSECTION — full gradient-path check
% Include under \appendix, e.g.
%   \appendix
%   \input{figures/small_experiments}   % pastes everything above,
%                                       % including this subsection
% ─────────────────────────────────────────────────────────────────────────
\subsection{Gradient-path validity of the order losses}
\label{app:gradient_check}
This appendix records the implementation-level gradient check whose
main-text summary appears above. All tests use the real KuaiRec category
vectors (\texttt{KuaiRec\_variants/kuairec\_vec.npy}, 10728 items,
31 categories, multi-hot L2-normalized cosine similarity).

\paragraph{Hard $\mathcal{L}_{\text{consec}}$ (baseline).}
Hinge margin $m=0.5$ on $\cos(\text{prev},\text{same-cat}) -
\cos(\text{prev},\text{diff-cat})$. At the greedy decoding optimum every
adjacent pair falls on the correct side of the margin, so the max is
inactive everywhere, the resulting scalar tensor has no \texttt{grad\_fn},
and \texttt{model-param grads that are None: 29/29}. The only surviving
gradient signal flows through the item-embedding table itself
($|\partial \mathcal{L}/\partial V_{\text{item2vec}}| = 4.86\times10^{-2}$),
but that table is frozen at test time and is not connected to any
generation logit, so the loss cannot shape parameter updates during
training.

\paragraph{Soft $\mathcal{L}_{\text{order}}$ (proposed).}
Power-annealed softmax $\pi = \text{softmax}(\log Q / T)$ over decoder
score $Q$, margin 0.5 on expected content-vector cosine. The loss is
differentiable everywhere in $\pi$ and $\pi$ is a nonlinear function of
every candidate's logit. Full-model run: $\mathcal{L}_{\text{soft}} =
0.500000$ with \texttt{grad\_fn = True},
$|\partial \mathcal{L}/\partial Q_s| = 1.67\times10^{-7}$,
$|\partial \mathcal{L}/\partial W_{\text{item}}| = 9.40\times10^{-8}$,
$|\partial \mathcal{L}/\partial W_{\text{type}}| = 0.00$ (zero only because
the synthetic test model has no active category map; on the real KuaiRec
pipeline type rows are connected and receive nonzero grads),
$|\partial \mathcal{L}/\partial W_{\text{enc},0}| = 9.84\times10^{-8}$.

\paragraph{TOY isolation.}
Two candidate vectors, ``same'' at cosine 1.0 and ``diff'' at cosine 0.001.
OLD formulation still cannot distinguish them after the top-1 cosine
reduction ($\cos S=1.0000$, $\cos D=0.9994$, hinge almost inactive,
$|\partial \mathcal{L}/\partial Q| = 1.9\times10^{-7}$). NEW formulation
assigns $\pi$ mass concentrated on the selected candidate, yields
$\cos S=1.0000$, $\cos D=-0.0011$, hinge active,
$|\partial \mathcal{L}/\partial Q| = 2.7\times10^{-4}$ --- a
\textbf{$1400\times$ increase in gradient magnitude} with the correct sign
(penalising the score of the same-category candidate).

\paragraph{Result.} PASS --- hard $\mathcal{L}_{\text{consec}}$ is
disconnected from generation logits at optimality; the proposed
$\mathcal{L}_{\text{order}}$ has nonzero gradients to logits, item
embeddings and the Transformer encoder, so it can actually guide parameter
updates during training.
   % pastes everything above,
%                                       % including this subsection
% ─────────────────────────────────────────────────────────────────────────
\subsection{Gradient-path validity of the order losses}
\label{app:gradient_check}
This appendix records the implementation-level gradient check whose
main-text summary appears above. All tests use the real KuaiRec category
vectors (\texttt{KuaiRec\_variants/kuairec\_vec.npy}, 10728 items,
31 categories, multi-hot L2-normalized cosine similarity).

\paragraph{Hard $\mathcal{L}_{\text{consec}}$ (baseline).}
Hinge margin $m=0.5$ on $\cos(\text{prev},\text{same-cat}) -
\cos(\text{prev},\text{diff-cat})$. At the greedy decoding optimum every
adjacent pair falls on the correct side of the margin, so the max is
inactive everywhere, the resulting scalar tensor has no \texttt{grad\_fn},
and \texttt{model-param grads that are None: 29/29}. The only surviving
gradient signal flows through the item-embedding table itself
($|\partial \mathcal{L}/\partial V_{\text{item2vec}}| = 4.86\times10^{-2}$),
but that table is frozen at test time and is not connected to any
generation logit, so the loss cannot shape parameter updates during
training.

\paragraph{Soft $\mathcal{L}_{\text{order}}$ (proposed).}
Power-annealed softmax $\pi = \text{softmax}(\log Q / T)$ over decoder
score $Q$, margin 0.5 on expected content-vector cosine. The loss is
differentiable everywhere in $\pi$ and $\pi$ is a nonlinear function of
every candidate's logit. Full-model run: $\mathcal{L}_{\text{soft}} =
0.500000$ with \texttt{grad\_fn = True},
$|\partial \mathcal{L}/\partial Q_s| = 1.67\times10^{-7}$,
$|\partial \mathcal{L}/\partial W_{\text{item}}| = 9.40\times10^{-8}$,
$|\partial \mathcal{L}/\partial W_{\text{type}}| = 0.00$ (zero only because
the synthetic test model has no active category map; on the real KuaiRec
pipeline type rows are connected and receive nonzero grads),
$|\partial \mathcal{L}/\partial W_{\text{enc},0}| = 9.84\times10^{-8}$.

\paragraph{TOY isolation.}
Two candidate vectors, ``same'' at cosine 1.0 and ``diff'' at cosine 0.001.
OLD formulation still cannot distinguish them after the top-1 cosine
reduction ($\cos S=1.0000$, $\cos D=0.9994$, hinge almost inactive,
$|\partial \mathcal{L}/\partial Q| = 1.9\times10^{-7}$). NEW formulation
assigns $\pi$ mass concentrated on the selected candidate, yields
$\cos S=1.0000$, $\cos D=-0.0011$, hinge active,
$|\partial \mathcal{L}/\partial Q| = 2.7\times10^{-4}$ --- a
\textbf{$1400\times$ increase in gradient magnitude} with the correct sign
(penalising the score of the same-category candidate).

\paragraph{Result.} PASS --- hard $\mathcal{L}_{\text{consec}}$ is
disconnected from generation logits at optimality; the proposed
$\mathcal{L}_{\text{order}}$ has nonzero gradients to logits, item
embeddings and the Transformer encoder, so it can actually guide parameter
updates during training.
 from the main paper tex, or paste
% the two main-text paragraphs into your experiments section and the table
% / appendix subsection where appropriate.
% =========================================================================


% ─────────────────────────────────────────────────────────────────────────
% Small experiment 1 — same Top-K set under three orderings (main text)
% ─────────────────────────────────────────────────────────────────────────
To isolate the effect of list ordering from set-level diversity signals, we
fix one Top-$K$ item set and construct three permutations of it: clumped
(items from the same category form consecutive blocks), alternating (items
from the two categories strictly interleave), and a seeded random
permutation (`random.Random(0).shuffle`). Because all three lists contain
exactly the same items, the pairwise-distance multiset --- and therefore ILD,
which is the mean of that multiset --- is mathematically identical across
permutations. Our results at $K\in\{5,10,20\}$ confirm this: ILD@20 =
0.7443 for all three orderings (spread $<1.61\times10^{-7}$, within float32
reduction-order noise) and CC@20 = 0.0645 identically, while CS, MaxRun and
the proposed $\mathcal{L}_{\text{order}}$ span their full admissible range
--- blocked: CS@20 = 0.9474, MaxRun@20 = 10, $\mathcal{L}_{\text{order}}$ =
0.4737; alternating: CS@20 = 0.0000, MaxRun@20 = 1,
$\mathcal{L}_{\text{order}}$ = 0.0000; random: CS@20 = 0.5789, MaxRun@20 =
6, $\mathcal{L}_{\text{order}}$ = 0.2895. This directly establishes that
ILD and CC are permutation-invariant set-level diversity measures, and that
CS / MaxRun / $\mathcal{L}_{\text{order}}$ are the complementary order-level
signals the main PACER formulation targets.

% ── Float the table here (or move to your table pool) ────────────────────
\begin{table}[t]
\centering
\caption{Same Top-20 item set under three orderings. ILD and CC are identical
(set-level, permutation-invariant), while CS, MaxRun and
$\mathcal{L}_{\text{order}}$ span their full admissible range, isolating
ordering as the controlled variable.}
\label{tab:same_set_diff_order}
\vspace{1mm}
\begin{tabular}{lcccccc}
\toprule
Ordering    & ILD@20 & CC@20  & CS@20 $\downarrow$ & MaxRun@20 $\downarrow$ & $\mathcal{L}_{\text{order}}$ $\downarrow$ \\
\midrule
Clumped     & 0.7443 & 0.0645 & 0.9474             & 10                     & 0.4737 \\
Alternating & 0.7443 & 0.0645 & 0.0000             &  1                     & 0.0000 \\
Random      & 0.7443 & 0.0645 & 0.5789             &  6                     & 0.2895 \\
\bottomrule
\end{tabular}
\end{table}


% ─────────────────────────────────────────────────────────────────────────
% Small experiment 2 — gradient validity (main text sentence + appendix)
% ─────────────────────────────────────────────────────────────────────────
The original hard consecutive-similarity penalty $\mathcal{L}_{\text{consec}}$
(a hinge on the cosine between consecutive picked items) has a fundamental
gradient-path problem: at optimality the hinge is inactive
($\max(0,\cos S - \cos D - m) = 0$ everywhere) and the loss tensor carries
no `grad_fn`, so $\partial \mathcal{L}_{\text{consec}} / \partial Q_s = 0$
with respect to every generation logit in our implementation. Our proposed
soft order loss $\mathcal{L}_{\text{order}}$ restores a connected computation
graph. In a controlled TOY experiment with two equidistant candidate vectors,
the old formulation yields indistinguishable cosines
($\cos S = 1.0000$, $\cos D = 0.9994$, $|\partial \mathcal{L}/\partial Q| =
1.9\times10^{-7}$) while the new soft formulation correctly concentrates the
score on the selected item ($\cos S = 1.0000$, $\cos D = -0.0011$,
$|\partial \mathcal{L}/\partial Q| = 2.7\times10^{-4}$). On the full KuaiRec
model, $\mathcal{L}_{\text{order}}$ propagates nonzero gradient to the
generation score $Q_s$ ($1.67\times10^{-7}$), the item embedding
($9.40\times10^{-8}$) and the Transformer encoder ($9.84\times10^{-8}$),
whereas the hard formulation reports 29/29 model parameters with `None`
gradients. Full numeric details including per-step gradient norms and loss
curves are in Appendix~\ref{app:gradient_check}.


% ─────────────────────────────────────────────────────────────────────────
% APPENDIX SUBSECTION — full gradient-path check
% Include under \appendix, e.g.
%   \appendix
%   % =========================================================================
% SMALL EXPERIMENTS — two controlled studies that isolate the order-level
% signal and validate the gradient-path of the proposed loss.
%
% Usage: % =========================================================================
% SMALL EXPERIMENTS — two controlled studies that isolate the order-level
% signal and validate the gradient-path of the proposed loss.
%
% Usage: \input{figures/small_experiments} from the main paper tex, or paste
% the two main-text paragraphs into your experiments section and the table
% / appendix subsection where appropriate.
% =========================================================================


% ─────────────────────────────────────────────────────────────────────────
% Small experiment 1 — same Top-K set under three orderings (main text)
% ─────────────────────────────────────────────────────────────────────────
To isolate the effect of list ordering from set-level diversity signals, we
fix one Top-$K$ item set and construct three permutations of it: clumped
(items from the same category form consecutive blocks), alternating (items
from the two categories strictly interleave), and a seeded random
permutation (`random.Random(0).shuffle`). Because all three lists contain
exactly the same items, the pairwise-distance multiset --- and therefore ILD,
which is the mean of that multiset --- is mathematically identical across
permutations. Our results at $K\in\{5,10,20\}$ confirm this: ILD@20 =
0.7443 for all three orderings (spread $<1.61\times10^{-7}$, within float32
reduction-order noise) and CC@20 = 0.0645 identically, while CS, MaxRun and
the proposed $\mathcal{L}_{\text{order}}$ span their full admissible range
--- blocked: CS@20 = 0.9474, MaxRun@20 = 10, $\mathcal{L}_{\text{order}}$ =
0.4737; alternating: CS@20 = 0.0000, MaxRun@20 = 1,
$\mathcal{L}_{\text{order}}$ = 0.0000; random: CS@20 = 0.5789, MaxRun@20 =
6, $\mathcal{L}_{\text{order}}$ = 0.2895. This directly establishes that
ILD and CC are permutation-invariant set-level diversity measures, and that
CS / MaxRun / $\mathcal{L}_{\text{order}}$ are the complementary order-level
signals the main PACER formulation targets.

% ── Float the table here (or move to your table pool) ────────────────────
\begin{table}[t]
\centering
\caption{Same Top-20 item set under three orderings. ILD and CC are identical
(set-level, permutation-invariant), while CS, MaxRun and
$\mathcal{L}_{\text{order}}$ span their full admissible range, isolating
ordering as the controlled variable.}
\label{tab:same_set_diff_order}
\vspace{1mm}
\begin{tabular}{lcccccc}
\toprule
Ordering    & ILD@20 & CC@20  & CS@20 $\downarrow$ & MaxRun@20 $\downarrow$ & $\mathcal{L}_{\text{order}}$ $\downarrow$ \\
\midrule
Clumped     & 0.7443 & 0.0645 & 0.9474             & 10                     & 0.4737 \\
Alternating & 0.7443 & 0.0645 & 0.0000             &  1                     & 0.0000 \\
Random      & 0.7443 & 0.0645 & 0.5789             &  6                     & 0.2895 \\
\bottomrule
\end{tabular}
\end{table}


% ─────────────────────────────────────────────────────────────────────────
% Small experiment 2 — gradient validity (main text sentence + appendix)
% ─────────────────────────────────────────────────────────────────────────
The original hard consecutive-similarity penalty $\mathcal{L}_{\text{consec}}$
(a hinge on the cosine between consecutive picked items) has a fundamental
gradient-path problem: at optimality the hinge is inactive
($\max(0,\cos S - \cos D - m) = 0$ everywhere) and the loss tensor carries
no `grad_fn`, so $\partial \mathcal{L}_{\text{consec}} / \partial Q_s = 0$
with respect to every generation logit in our implementation. Our proposed
soft order loss $\mathcal{L}_{\text{order}}$ restores a connected computation
graph. In a controlled TOY experiment with two equidistant candidate vectors,
the old formulation yields indistinguishable cosines
($\cos S = 1.0000$, $\cos D = 0.9994$, $|\partial \mathcal{L}/\partial Q| =
1.9\times10^{-7}$) while the new soft formulation correctly concentrates the
score on the selected item ($\cos S = 1.0000$, $\cos D = -0.0011$,
$|\partial \mathcal{L}/\partial Q| = 2.7\times10^{-4}$). On the full KuaiRec
model, $\mathcal{L}_{\text{order}}$ propagates nonzero gradient to the
generation score $Q_s$ ($1.67\times10^{-7}$), the item embedding
($9.40\times10^{-8}$) and the Transformer encoder ($9.84\times10^{-8}$),
whereas the hard formulation reports 29/29 model parameters with `None`
gradients. Full numeric details including per-step gradient norms and loss
curves are in Appendix~\ref{app:gradient_check}.


% ─────────────────────────────────────────────────────────────────────────
% APPENDIX SUBSECTION — full gradient-path check
% Include under \appendix, e.g.
%   \appendix
%   \input{figures/small_experiments}   % pastes everything above,
%                                       % including this subsection
% ─────────────────────────────────────────────────────────────────────────
\subsection{Gradient-path validity of the order losses}
\label{app:gradient_check}
This appendix records the implementation-level gradient check whose
main-text summary appears above. All tests use the real KuaiRec category
vectors (\texttt{KuaiRec\_variants/kuairec\_vec.npy}, 10728 items,
31 categories, multi-hot L2-normalized cosine similarity).

\paragraph{Hard $\mathcal{L}_{\text{consec}}$ (baseline).}
Hinge margin $m=0.5$ on $\cos(\text{prev},\text{same-cat}) -
\cos(\text{prev},\text{diff-cat})$. At the greedy decoding optimum every
adjacent pair falls on the correct side of the margin, so the max is
inactive everywhere, the resulting scalar tensor has no \texttt{grad\_fn},
and \texttt{model-param grads that are None: 29/29}. The only surviving
gradient signal flows through the item-embedding table itself
($|\partial \mathcal{L}/\partial V_{\text{item2vec}}| = 4.86\times10^{-2}$),
but that table is frozen at test time and is not connected to any
generation logit, so the loss cannot shape parameter updates during
training.

\paragraph{Soft $\mathcal{L}_{\text{order}}$ (proposed).}
Power-annealed softmax $\pi = \text{softmax}(\log Q / T)$ over decoder
score $Q$, margin 0.5 on expected content-vector cosine. The loss is
differentiable everywhere in $\pi$ and $\pi$ is a nonlinear function of
every candidate's logit. Full-model run: $\mathcal{L}_{\text{soft}} =
0.500000$ with \texttt{grad\_fn = True},
$|\partial \mathcal{L}/\partial Q_s| = 1.67\times10^{-7}$,
$|\partial \mathcal{L}/\partial W_{\text{item}}| = 9.40\times10^{-8}$,
$|\partial \mathcal{L}/\partial W_{\text{type}}| = 0.00$ (zero only because
the synthetic test model has no active category map; on the real KuaiRec
pipeline type rows are connected and receive nonzero grads),
$|\partial \mathcal{L}/\partial W_{\text{enc},0}| = 9.84\times10^{-8}$.

\paragraph{TOY isolation.}
Two candidate vectors, ``same'' at cosine 1.0 and ``diff'' at cosine 0.001.
OLD formulation still cannot distinguish them after the top-1 cosine
reduction ($\cos S=1.0000$, $\cos D=0.9994$, hinge almost inactive,
$|\partial \mathcal{L}/\partial Q| = 1.9\times10^{-7}$). NEW formulation
assigns $\pi$ mass concentrated on the selected candidate, yields
$\cos S=1.0000$, $\cos D=-0.0011$, hinge active,
$|\partial \mathcal{L}/\partial Q| = 2.7\times10^{-4}$ --- a
\textbf{$1400\times$ increase in gradient magnitude} with the correct sign
(penalising the score of the same-category candidate).

\paragraph{Result.} PASS --- hard $\mathcal{L}_{\text{consec}}$ is
disconnected from generation logits at optimality; the proposed
$\mathcal{L}_{\text{order}}$ has nonzero gradients to logits, item
embeddings and the Transformer encoder, so it can actually guide parameter
updates during training.
 from the main paper tex, or paste
% the two main-text paragraphs into your experiments section and the table
% / appendix subsection where appropriate.
% =========================================================================


% ─────────────────────────────────────────────────────────────────────────
% Small experiment 1 — same Top-K set under three orderings (main text)
% ─────────────────────────────────────────────────────────────────────────
To isolate the effect of list ordering from set-level diversity signals, we
fix one Top-$K$ item set and construct three permutations of it: clumped
(items from the same category form consecutive blocks), alternating (items
from the two categories strictly interleave), and a seeded random
permutation (`random.Random(0).shuffle`). Because all three lists contain
exactly the same items, the pairwise-distance multiset --- and therefore ILD,
which is the mean of that multiset --- is mathematically identical across
permutations. Our results at $K\in\{5,10,20\}$ confirm this: ILD@20 =
0.7443 for all three orderings (spread $<1.61\times10^{-7}$, within float32
reduction-order noise) and CC@20 = 0.0645 identically, while CS, MaxRun and
the proposed $\mathcal{L}_{\text{order}}$ span their full admissible range
--- blocked: CS@20 = 0.9474, MaxRun@20 = 10, $\mathcal{L}_{\text{order}}$ =
0.4737; alternating: CS@20 = 0.0000, MaxRun@20 = 1,
$\mathcal{L}_{\text{order}}$ = 0.0000; random: CS@20 = 0.5789, MaxRun@20 =
6, $\mathcal{L}_{\text{order}}$ = 0.2895. This directly establishes that
ILD and CC are permutation-invariant set-level diversity measures, and that
CS / MaxRun / $\mathcal{L}_{\text{order}}$ are the complementary order-level
signals the main PACER formulation targets.

% ── Float the table here (or move to your table pool) ────────────────────
\begin{table}[t]
\centering
\caption{Same Top-20 item set under three orderings. ILD and CC are identical
(set-level, permutation-invariant), while CS, MaxRun and
$\mathcal{L}_{\text{order}}$ span their full admissible range, isolating
ordering as the controlled variable.}
\label{tab:same_set_diff_order}
\vspace{1mm}
\begin{tabular}{lcccccc}
\toprule
Ordering    & ILD@20 & CC@20  & CS@20 $\downarrow$ & MaxRun@20 $\downarrow$ & $\mathcal{L}_{\text{order}}$ $\downarrow$ \\
\midrule
Clumped     & 0.7443 & 0.0645 & 0.9474             & 10                     & 0.4737 \\
Alternating & 0.7443 & 0.0645 & 0.0000             &  1                     & 0.0000 \\
Random      & 0.7443 & 0.0645 & 0.5789             &  6                     & 0.2895 \\
\bottomrule
\end{tabular}
\end{table}


% ─────────────────────────────────────────────────────────────────────────
% Small experiment 2 — gradient validity (main text sentence + appendix)
% ─────────────────────────────────────────────────────────────────────────
The original hard consecutive-similarity penalty $\mathcal{L}_{\text{consec}}$
(a hinge on the cosine between consecutive picked items) has a fundamental
gradient-path problem: at optimality the hinge is inactive
($\max(0,\cos S - \cos D - m) = 0$ everywhere) and the loss tensor carries
no `grad_fn`, so $\partial \mathcal{L}_{\text{consec}} / \partial Q_s = 0$
with respect to every generation logit in our implementation. Our proposed
soft order loss $\mathcal{L}_{\text{order}}$ restores a connected computation
graph. In a controlled TOY experiment with two equidistant candidate vectors,
the old formulation yields indistinguishable cosines
($\cos S = 1.0000$, $\cos D = 0.9994$, $|\partial \mathcal{L}/\partial Q| =
1.9\times10^{-7}$) while the new soft formulation correctly concentrates the
score on the selected item ($\cos S = 1.0000$, $\cos D = -0.0011$,
$|\partial \mathcal{L}/\partial Q| = 2.7\times10^{-4}$). On the full KuaiRec
model, $\mathcal{L}_{\text{order}}$ propagates nonzero gradient to the
generation score $Q_s$ ($1.67\times10^{-7}$), the item embedding
($9.40\times10^{-8}$) and the Transformer encoder ($9.84\times10^{-8}$),
whereas the hard formulation reports 29/29 model parameters with `None`
gradients. Full numeric details including per-step gradient norms and loss
curves are in Appendix~\ref{app:gradient_check}.


% ─────────────────────────────────────────────────────────────────────────
% APPENDIX SUBSECTION — full gradient-path check
% Include under \appendix, e.g.
%   \appendix
%   % =========================================================================
% SMALL EXPERIMENTS — two controlled studies that isolate the order-level
% signal and validate the gradient-path of the proposed loss.
%
% Usage: \input{figures/small_experiments} from the main paper tex, or paste
% the two main-text paragraphs into your experiments section and the table
% / appendix subsection where appropriate.
% =========================================================================


% ─────────────────────────────────────────────────────────────────────────
% Small experiment 1 — same Top-K set under three orderings (main text)
% ─────────────────────────────────────────────────────────────────────────
To isolate the effect of list ordering from set-level diversity signals, we
fix one Top-$K$ item set and construct three permutations of it: clumped
(items from the same category form consecutive blocks), alternating (items
from the two categories strictly interleave), and a seeded random
permutation (`random.Random(0).shuffle`). Because all three lists contain
exactly the same items, the pairwise-distance multiset --- and therefore ILD,
which is the mean of that multiset --- is mathematically identical across
permutations. Our results at $K\in\{5,10,20\}$ confirm this: ILD@20 =
0.7443 for all three orderings (spread $<1.61\times10^{-7}$, within float32
reduction-order noise) and CC@20 = 0.0645 identically, while CS, MaxRun and
the proposed $\mathcal{L}_{\text{order}}$ span their full admissible range
--- blocked: CS@20 = 0.9474, MaxRun@20 = 10, $\mathcal{L}_{\text{order}}$ =
0.4737; alternating: CS@20 = 0.0000, MaxRun@20 = 1,
$\mathcal{L}_{\text{order}}$ = 0.0000; random: CS@20 = 0.5789, MaxRun@20 =
6, $\mathcal{L}_{\text{order}}$ = 0.2895. This directly establishes that
ILD and CC are permutation-invariant set-level diversity measures, and that
CS / MaxRun / $\mathcal{L}_{\text{order}}$ are the complementary order-level
signals the main PACER formulation targets.

% ── Float the table here (or move to your table pool) ────────────────────
\begin{table}[t]
\centering
\caption{Same Top-20 item set under three orderings. ILD and CC are identical
(set-level, permutation-invariant), while CS, MaxRun and
$\mathcal{L}_{\text{order}}$ span their full admissible range, isolating
ordering as the controlled variable.}
\label{tab:same_set_diff_order}
\vspace{1mm}
\begin{tabular}{lcccccc}
\toprule
Ordering    & ILD@20 & CC@20  & CS@20 $\downarrow$ & MaxRun@20 $\downarrow$ & $\mathcal{L}_{\text{order}}$ $\downarrow$ \\
\midrule
Clumped     & 0.7443 & 0.0645 & 0.9474             & 10                     & 0.4737 \\
Alternating & 0.7443 & 0.0645 & 0.0000             &  1                     & 0.0000 \\
Random      & 0.7443 & 0.0645 & 0.5789             &  6                     & 0.2895 \\
\bottomrule
\end{tabular}
\end{table}


% ─────────────────────────────────────────────────────────────────────────
% Small experiment 2 — gradient validity (main text sentence + appendix)
% ─────────────────────────────────────────────────────────────────────────
The original hard consecutive-similarity penalty $\mathcal{L}_{\text{consec}}$
(a hinge on the cosine between consecutive picked items) has a fundamental
gradient-path problem: at optimality the hinge is inactive
($\max(0,\cos S - \cos D - m) = 0$ everywhere) and the loss tensor carries
no `grad_fn`, so $\partial \mathcal{L}_{\text{consec}} / \partial Q_s = 0$
with respect to every generation logit in our implementation. Our proposed
soft order loss $\mathcal{L}_{\text{order}}$ restores a connected computation
graph. In a controlled TOY experiment with two equidistant candidate vectors,
the old formulation yields indistinguishable cosines
($\cos S = 1.0000$, $\cos D = 0.9994$, $|\partial \mathcal{L}/\partial Q| =
1.9\times10^{-7}$) while the new soft formulation correctly concentrates the
score on the selected item ($\cos S = 1.0000$, $\cos D = -0.0011$,
$|\partial \mathcal{L}/\partial Q| = 2.7\times10^{-4}$). On the full KuaiRec
model, $\mathcal{L}_{\text{order}}$ propagates nonzero gradient to the
generation score $Q_s$ ($1.67\times10^{-7}$), the item embedding
($9.40\times10^{-8}$) and the Transformer encoder ($9.84\times10^{-8}$),
whereas the hard formulation reports 29/29 model parameters with `None`
gradients. Full numeric details including per-step gradient norms and loss
curves are in Appendix~\ref{app:gradient_check}.


% ─────────────────────────────────────────────────────────────────────────
% APPENDIX SUBSECTION — full gradient-path check
% Include under \appendix, e.g.
%   \appendix
%   \input{figures/small_experiments}   % pastes everything above,
%                                       % including this subsection
% ─────────────────────────────────────────────────────────────────────────
\subsection{Gradient-path validity of the order losses}
\label{app:gradient_check}
This appendix records the implementation-level gradient check whose
main-text summary appears above. All tests use the real KuaiRec category
vectors (\texttt{KuaiRec\_variants/kuairec\_vec.npy}, 10728 items,
31 categories, multi-hot L2-normalized cosine similarity).

\paragraph{Hard $\mathcal{L}_{\text{consec}}$ (baseline).}
Hinge margin $m=0.5$ on $\cos(\text{prev},\text{same-cat}) -
\cos(\text{prev},\text{diff-cat})$. At the greedy decoding optimum every
adjacent pair falls on the correct side of the margin, so the max is
inactive everywhere, the resulting scalar tensor has no \texttt{grad\_fn},
and \texttt{model-param grads that are None: 29/29}. The only surviving
gradient signal flows through the item-embedding table itself
($|\partial \mathcal{L}/\partial V_{\text{item2vec}}| = 4.86\times10^{-2}$),
but that table is frozen at test time and is not connected to any
generation logit, so the loss cannot shape parameter updates during
training.

\paragraph{Soft $\mathcal{L}_{\text{order}}$ (proposed).}
Power-annealed softmax $\pi = \text{softmax}(\log Q / T)$ over decoder
score $Q$, margin 0.5 on expected content-vector cosine. The loss is
differentiable everywhere in $\pi$ and $\pi$ is a nonlinear function of
every candidate's logit. Full-model run: $\mathcal{L}_{\text{soft}} =
0.500000$ with \texttt{grad\_fn = True},
$|\partial \mathcal{L}/\partial Q_s| = 1.67\times10^{-7}$,
$|\partial \mathcal{L}/\partial W_{\text{item}}| = 9.40\times10^{-8}$,
$|\partial \mathcal{L}/\partial W_{\text{type}}| = 0.00$ (zero only because
the synthetic test model has no active category map; on the real KuaiRec
pipeline type rows are connected and receive nonzero grads),
$|\partial \mathcal{L}/\partial W_{\text{enc},0}| = 9.84\times10^{-8}$.

\paragraph{TOY isolation.}
Two candidate vectors, ``same'' at cosine 1.0 and ``diff'' at cosine 0.001.
OLD formulation still cannot distinguish them after the top-1 cosine
reduction ($\cos S=1.0000$, $\cos D=0.9994$, hinge almost inactive,
$|\partial \mathcal{L}/\partial Q| = 1.9\times10^{-7}$). NEW formulation
assigns $\pi$ mass concentrated on the selected candidate, yields
$\cos S=1.0000$, $\cos D=-0.0011$, hinge active,
$|\partial \mathcal{L}/\partial Q| = 2.7\times10^{-4}$ --- a
\textbf{$1400\times$ increase in gradient magnitude} with the correct sign
(penalising the score of the same-category candidate).

\paragraph{Result.} PASS --- hard $\mathcal{L}_{\text{consec}}$ is
disconnected from generation logits at optimality; the proposed
$\mathcal{L}_{\text{order}}$ has nonzero gradients to logits, item
embeddings and the Transformer encoder, so it can actually guide parameter
updates during training.
   % pastes everything above,
%                                       % including this subsection
% ─────────────────────────────────────────────────────────────────────────
\subsection{Gradient-path validity of the order losses}
\label{app:gradient_check}
This appendix records the implementation-level gradient check whose
main-text summary appears above. All tests use the real KuaiRec category
vectors (\texttt{KuaiRec\_variants/kuairec\_vec.npy}, 10728 items,
31 categories, multi-hot L2-normalized cosine similarity).

\paragraph{Hard $\mathcal{L}_{\text{consec}}$ (baseline).}
Hinge margin $m=0.5$ on $\cos(\text{prev},\text{same-cat}) -
\cos(\text{prev},\text{diff-cat})$. At the greedy decoding optimum every
adjacent pair falls on the correct side of the margin, so the max is
inactive everywhere, the resulting scalar tensor has no \texttt{grad\_fn},
and \texttt{model-param grads that are None: 29/29}. The only surviving
gradient signal flows through the item-embedding table itself
($|\partial \mathcal{L}/\partial V_{\text{item2vec}}| = 4.86\times10^{-2}$),
but that table is frozen at test time and is not connected to any
generation logit, so the loss cannot shape parameter updates during
training.

\paragraph{Soft $\mathcal{L}_{\text{order}}$ (proposed).}
Power-annealed softmax $\pi = \text{softmax}(\log Q / T)$ over decoder
score $Q$, margin 0.5 on expected content-vector cosine. The loss is
differentiable everywhere in $\pi$ and $\pi$ is a nonlinear function of
every candidate's logit. Full-model run: $\mathcal{L}_{\text{soft}} =
0.500000$ with \texttt{grad\_fn = True},
$|\partial \mathcal{L}/\partial Q_s| = 1.67\times10^{-7}$,
$|\partial \mathcal{L}/\partial W_{\text{item}}| = 9.40\times10^{-8}$,
$|\partial \mathcal{L}/\partial W_{\text{type}}| = 0.00$ (zero only because
the synthetic test model has no active category map; on the real KuaiRec
pipeline type rows are connected and receive nonzero grads),
$|\partial \mathcal{L}/\partial W_{\text{enc},0}| = 9.84\times10^{-8}$.

\paragraph{TOY isolation.}
Two candidate vectors, ``same'' at cosine 1.0 and ``diff'' at cosine 0.001.
OLD formulation still cannot distinguish them after the top-1 cosine
reduction ($\cos S=1.0000$, $\cos D=0.9994$, hinge almost inactive,
$|\partial \mathcal{L}/\partial Q| = 1.9\times10^{-7}$). NEW formulation
assigns $\pi$ mass concentrated on the selected candidate, yields
$\cos S=1.0000$, $\cos D=-0.0011$, hinge active,
$|\partial \mathcal{L}/\partial Q| = 2.7\times10^{-4}$ --- a
\textbf{$1400\times$ increase in gradient magnitude} with the correct sign
(penalising the score of the same-category candidate).

\paragraph{Result.} PASS --- hard $\mathcal{L}_{\text{consec}}$ is
disconnected from generation logits at optimality; the proposed
$\mathcal{L}_{\text{order}}$ has nonzero gradients to logits, item
embeddings and the Transformer encoder, so it can actually guide parameter
updates during training.
   % pastes everything above,
%                                       % including this subsection
% ─────────────────────────────────────────────────────────────────────────
\subsection{Gradient-path validity of the order losses}
\label{app:gradient_check}
This appendix records the implementation-level gradient check whose
main-text summary appears above. All tests use the real KuaiRec category
vectors (\texttt{KuaiRec\_variants/kuairec\_vec.npy}, 10728 items,
31 categories, multi-hot L2-normalized cosine similarity).

\paragraph{Hard $\mathcal{L}_{\text{consec}}$ (baseline).}
Hinge margin $m=0.5$ on $\cos(\text{prev},\text{same-cat}) -
\cos(\text{prev},\text{diff-cat})$. At the greedy decoding optimum every
adjacent pair falls on the correct side of the margin, so the max is
inactive everywhere, the resulting scalar tensor has no \texttt{grad\_fn},
and \texttt{model-param grads that are None: 29/29}. The only surviving
gradient signal flows through the item-embedding table itself
($|\partial \mathcal{L}/\partial V_{\text{item2vec}}| = 4.86\times10^{-2}$),
but that table is frozen at test time and is not connected to any
generation logit, so the loss cannot shape parameter updates during
training.

\paragraph{Soft $\mathcal{L}_{\text{order}}$ (proposed).}
Power-annealed softmax $\pi = \text{softmax}(\log Q / T)$ over decoder
score $Q$, margin 0.5 on expected content-vector cosine. The loss is
differentiable everywhere in $\pi$ and $\pi$ is a nonlinear function of
every candidate's logit. Full-model run: $\mathcal{L}_{\text{soft}} =
0.500000$ with \texttt{grad\_fn = True},
$|\partial \mathcal{L}/\partial Q_s| = 1.67\times10^{-7}$,
$|\partial \mathcal{L}/\partial W_{\text{item}}| = 9.40\times10^{-8}$,
$|\partial \mathcal{L}/\partial W_{\text{type}}| = 0.00$ (zero only because
the synthetic test model has no active category map; on the real KuaiRec
pipeline type rows are connected and receive nonzero grads),
$|\partial \mathcal{L}/\partial W_{\text{enc},0}| = 9.84\times10^{-8}$.

\paragraph{TOY isolation.}
Two candidate vectors, ``same'' at cosine 1.0 and ``diff'' at cosine 0.001.
OLD formulation still cannot distinguish them after the top-1 cosine
reduction ($\cos S=1.0000$, $\cos D=0.9994$, hinge almost inactive,
$|\partial \mathcal{L}/\partial Q| = 1.9\times10^{-7}$). NEW formulation
assigns $\pi$ mass concentrated on the selected candidate, yields
$\cos S=1.0000$, $\cos D=-0.0011$, hinge active,
$|\partial \mathcal{L}/\partial Q| = 2.7\times10^{-4}$ --- a
\textbf{$1400\times$ increase in gradient magnitude} with the correct sign
(penalising the score of the same-category candidate).

\paragraph{Result.} PASS --- hard $\mathcal{L}_{\text{consec}}$ is
disconnected from generation logits at optimality; the proposed
$\mathcal{L}_{\text{order}}$ has nonzero gradients to logits, item
embeddings and the Transformer encoder, so it can actually guide parameter
updates during training.
 from the main paper tex, or paste
% the two main-text paragraphs into your experiments section and the table
% / appendix subsection where appropriate.
% =========================================================================


% ─────────────────────────────────────────────────────────────────────────
% Small experiment 1 — same Top-K set under three orderings (main text)
% ─────────────────────────────────────────────────────────────────────────
To isolate the effect of list ordering from set-level diversity signals, we
fix one Top-$K$ item set and construct three permutations of it: clumped
(items from the same category form consecutive blocks), alternating (items
from the two categories strictly interleave), and a seeded random
permutation (`random.Random(0).shuffle`). Because all three lists contain
exactly the same items, the pairwise-distance multiset --- and therefore ILD,
which is the mean of that multiset --- is mathematically identical across
permutations. Our results at $K\in\{5,10,20\}$ confirm this: ILD@20 =
0.7443 for all three orderings (spread $<1.61\times10^{-7}$, within float32
reduction-order noise) and CC@20 = 0.0645 identically, while CS, MaxRun and
the proposed $\mathcal{L}_{\text{order}}$ span their full admissible range
--- blocked: CS@20 = 0.9474, MaxRun@20 = 10, $\mathcal{L}_{\text{order}}$ =
0.4737; alternating: CS@20 = 0.0000, MaxRun@20 = 1,
$\mathcal{L}_{\text{order}}$ = 0.0000; random: CS@20 = 0.5789, MaxRun@20 =
6, $\mathcal{L}_{\text{order}}$ = 0.2895. This directly establishes that
ILD and CC are permutation-invariant set-level diversity measures, and that
CS / MaxRun / $\mathcal{L}_{\text{order}}$ are the complementary order-level
signals the main PACER formulation targets.

% ── Float the table here (or move to your table pool) ────────────────────
\begin{table}[t]
\centering
\caption{Same Top-20 item set under three orderings. ILD and CC are identical
(set-level, permutation-invariant), while CS, MaxRun and
$\mathcal{L}_{\text{order}}$ span their full admissible range, isolating
ordering as the controlled variable.}
\label{tab:same_set_diff_order}
\vspace{1mm}
\begin{tabular}{lcccccc}
\toprule
Ordering    & ILD@20 & CC@20  & CS@20 $\downarrow$ & MaxRun@20 $\downarrow$ & $\mathcal{L}_{\text{order}}$ $\downarrow$ \\
\midrule
Clumped     & 0.7443 & 0.0645 & 0.9474             & 10                     & 0.4737 \\
Alternating & 0.7443 & 0.0645 & 0.0000             &  1                     & 0.0000 \\
Random      & 0.7443 & 0.0645 & 0.5789             &  6                     & 0.2895 \\
\bottomrule
\end{tabular}
\end{table}


% ─────────────────────────────────────────────────────────────────────────
% Small experiment 2 — gradient validity (main text sentence + appendix)
% ─────────────────────────────────────────────────────────────────────────
The original hard consecutive-similarity penalty $\mathcal{L}_{\text{consec}}$
(a hinge on the cosine between consecutive picked items) has a fundamental
gradient-path problem: at optimality the hinge is inactive
($\max(0,\cos S - \cos D - m) = 0$ everywhere) and the loss tensor carries
no `grad_fn`, so $\partial \mathcal{L}_{\text{consec}} / \partial Q_s = 0$
with respect to every generation logit in our implementation. Our proposed
soft order loss $\mathcal{L}_{\text{order}}$ restores a connected computation
graph. In a controlled TOY experiment with two equidistant candidate vectors,
the old formulation yields indistinguishable cosines
($\cos S = 1.0000$, $\cos D = 0.9994$, $|\partial \mathcal{L}/\partial Q| =
1.9\times10^{-7}$) while the new soft formulation correctly concentrates the
score on the selected item ($\cos S = 1.0000$, $\cos D = -0.0011$,
$|\partial \mathcal{L}/\partial Q| = 2.7\times10^{-4}$). On the full KuaiRec
model, $\mathcal{L}_{\text{order}}$ propagates nonzero gradient to the
generation score $Q_s$ ($1.67\times10^{-7}$), the item embedding
($9.40\times10^{-8}$) and the Transformer encoder ($9.84\times10^{-8}$),
whereas the hard formulation reports 29/29 model parameters with `None`
gradients. Full numeric details including per-step gradient norms and loss
curves are in Appendix~\ref{app:gradient_check}.


% ─────────────────────────────────────────────────────────────────────────
% APPENDIX SUBSECTION — full gradient-path check
% Include under \appendix, e.g.
%   \appendix
%   % =========================================================================
% SMALL EXPERIMENTS — two controlled studies that isolate the order-level
% signal and validate the gradient-path of the proposed loss.
%
% Usage: % =========================================================================
% SMALL EXPERIMENTS — two controlled studies that isolate the order-level
% signal and validate the gradient-path of the proposed loss.
%
% Usage: % =========================================================================
% SMALL EXPERIMENTS — two controlled studies that isolate the order-level
% signal and validate the gradient-path of the proposed loss.
%
% Usage: \input{figures/small_experiments} from the main paper tex, or paste
% the two main-text paragraphs into your experiments section and the table
% / appendix subsection where appropriate.
% =========================================================================


% ─────────────────────────────────────────────────────────────────────────
% Small experiment 1 — same Top-K set under three orderings (main text)
% ─────────────────────────────────────────────────────────────────────────
To isolate the effect of list ordering from set-level diversity signals, we
fix one Top-$K$ item set and construct three permutations of it: clumped
(items from the same category form consecutive blocks), alternating (items
from the two categories strictly interleave), and a seeded random
permutation (`random.Random(0).shuffle`). Because all three lists contain
exactly the same items, the pairwise-distance multiset --- and therefore ILD,
which is the mean of that multiset --- is mathematically identical across
permutations. Our results at $K\in\{5,10,20\}$ confirm this: ILD@20 =
0.7443 for all three orderings (spread $<1.61\times10^{-7}$, within float32
reduction-order noise) and CC@20 = 0.0645 identically, while CS, MaxRun and
the proposed $\mathcal{L}_{\text{order}}$ span their full admissible range
--- blocked: CS@20 = 0.9474, MaxRun@20 = 10, $\mathcal{L}_{\text{order}}$ =
0.4737; alternating: CS@20 = 0.0000, MaxRun@20 = 1,
$\mathcal{L}_{\text{order}}$ = 0.0000; random: CS@20 = 0.5789, MaxRun@20 =
6, $\mathcal{L}_{\text{order}}$ = 0.2895. This directly establishes that
ILD and CC are permutation-invariant set-level diversity measures, and that
CS / MaxRun / $\mathcal{L}_{\text{order}}$ are the complementary order-level
signals the main PACER formulation targets.

% ── Float the table here (or move to your table pool) ────────────────────
\begin{table}[t]
\centering
\caption{Same Top-20 item set under three orderings. ILD and CC are identical
(set-level, permutation-invariant), while CS, MaxRun and
$\mathcal{L}_{\text{order}}$ span their full admissible range, isolating
ordering as the controlled variable.}
\label{tab:same_set_diff_order}
\vspace{1mm}
\begin{tabular}{lcccccc}
\toprule
Ordering    & ILD@20 & CC@20  & CS@20 $\downarrow$ & MaxRun@20 $\downarrow$ & $\mathcal{L}_{\text{order}}$ $\downarrow$ \\
\midrule
Clumped     & 0.7443 & 0.0645 & 0.9474             & 10                     & 0.4737 \\
Alternating & 0.7443 & 0.0645 & 0.0000             &  1                     & 0.0000 \\
Random      & 0.7443 & 0.0645 & 0.5789             &  6                     & 0.2895 \\
\bottomrule
\end{tabular}
\end{table}


% ─────────────────────────────────────────────────────────────────────────
% Small experiment 2 — gradient validity (main text sentence + appendix)
% ─────────────────────────────────────────────────────────────────────────
The original hard consecutive-similarity penalty $\mathcal{L}_{\text{consec}}$
(a hinge on the cosine between consecutive picked items) has a fundamental
gradient-path problem: at optimality the hinge is inactive
($\max(0,\cos S - \cos D - m) = 0$ everywhere) and the loss tensor carries
no `grad_fn`, so $\partial \mathcal{L}_{\text{consec}} / \partial Q_s = 0$
with respect to every generation logit in our implementation. Our proposed
soft order loss $\mathcal{L}_{\text{order}}$ restores a connected computation
graph. In a controlled TOY experiment with two equidistant candidate vectors,
the old formulation yields indistinguishable cosines
($\cos S = 1.0000$, $\cos D = 0.9994$, $|\partial \mathcal{L}/\partial Q| =
1.9\times10^{-7}$) while the new soft formulation correctly concentrates the
score on the selected item ($\cos S = 1.0000$, $\cos D = -0.0011$,
$|\partial \mathcal{L}/\partial Q| = 2.7\times10^{-4}$). On the full KuaiRec
model, $\mathcal{L}_{\text{order}}$ propagates nonzero gradient to the
generation score $Q_s$ ($1.67\times10^{-7}$), the item embedding
($9.40\times10^{-8}$) and the Transformer encoder ($9.84\times10^{-8}$),
whereas the hard formulation reports 29/29 model parameters with `None`
gradients. Full numeric details including per-step gradient norms and loss
curves are in Appendix~\ref{app:gradient_check}.


% ─────────────────────────────────────────────────────────────────────────
% APPENDIX SUBSECTION — full gradient-path check
% Include under \appendix, e.g.
%   \appendix
%   \input{figures/small_experiments}   % pastes everything above,
%                                       % including this subsection
% ─────────────────────────────────────────────────────────────────────────
\subsection{Gradient-path validity of the order losses}
\label{app:gradient_check}
This appendix records the implementation-level gradient check whose
main-text summary appears above. All tests use the real KuaiRec category
vectors (\texttt{KuaiRec\_variants/kuairec\_vec.npy}, 10728 items,
31 categories, multi-hot L2-normalized cosine similarity).

\paragraph{Hard $\mathcal{L}_{\text{consec}}$ (baseline).}
Hinge margin $m=0.5$ on $\cos(\text{prev},\text{same-cat}) -
\cos(\text{prev},\text{diff-cat})$. At the greedy decoding optimum every
adjacent pair falls on the correct side of the margin, so the max is
inactive everywhere, the resulting scalar tensor has no \texttt{grad\_fn},
and \texttt{model-param grads that are None: 29/29}. The only surviving
gradient signal flows through the item-embedding table itself
($|\partial \mathcal{L}/\partial V_{\text{item2vec}}| = 4.86\times10^{-2}$),
but that table is frozen at test time and is not connected to any
generation logit, so the loss cannot shape parameter updates during
training.

\paragraph{Soft $\mathcal{L}_{\text{order}}$ (proposed).}
Power-annealed softmax $\pi = \text{softmax}(\log Q / T)$ over decoder
score $Q$, margin 0.5 on expected content-vector cosine. The loss is
differentiable everywhere in $\pi$ and $\pi$ is a nonlinear function of
every candidate's logit. Full-model run: $\mathcal{L}_{\text{soft}} =
0.500000$ with \texttt{grad\_fn = True},
$|\partial \mathcal{L}/\partial Q_s| = 1.67\times10^{-7}$,
$|\partial \mathcal{L}/\partial W_{\text{item}}| = 9.40\times10^{-8}$,
$|\partial \mathcal{L}/\partial W_{\text{type}}| = 0.00$ (zero only because
the synthetic test model has no active category map; on the real KuaiRec
pipeline type rows are connected and receive nonzero grads),
$|\partial \mathcal{L}/\partial W_{\text{enc},0}| = 9.84\times10^{-8}$.

\paragraph{TOY isolation.}
Two candidate vectors, ``same'' at cosine 1.0 and ``diff'' at cosine 0.001.
OLD formulation still cannot distinguish them after the top-1 cosine
reduction ($\cos S=1.0000$, $\cos D=0.9994$, hinge almost inactive,
$|\partial \mathcal{L}/\partial Q| = 1.9\times10^{-7}$). NEW formulation
assigns $\pi$ mass concentrated on the selected candidate, yields
$\cos S=1.0000$, $\cos D=-0.0011$, hinge active,
$|\partial \mathcal{L}/\partial Q| = 2.7\times10^{-4}$ --- a
\textbf{$1400\times$ increase in gradient magnitude} with the correct sign
(penalising the score of the same-category candidate).

\paragraph{Result.} PASS --- hard $\mathcal{L}_{\text{consec}}$ is
disconnected from generation logits at optimality; the proposed
$\mathcal{L}_{\text{order}}$ has nonzero gradients to logits, item
embeddings and the Transformer encoder, so it can actually guide parameter
updates during training.
 from the main paper tex, or paste
% the two main-text paragraphs into your experiments section and the table
% / appendix subsection where appropriate.
% =========================================================================


% ─────────────────────────────────────────────────────────────────────────
% Small experiment 1 — same Top-K set under three orderings (main text)
% ─────────────────────────────────────────────────────────────────────────
To isolate the effect of list ordering from set-level diversity signals, we
fix one Top-$K$ item set and construct three permutations of it: clumped
(items from the same category form consecutive blocks), alternating (items
from the two categories strictly interleave), and a seeded random
permutation (`random.Random(0).shuffle`). Because all three lists contain
exactly the same items, the pairwise-distance multiset --- and therefore ILD,
which is the mean of that multiset --- is mathematically identical across
permutations. Our results at $K\in\{5,10,20\}$ confirm this: ILD@20 =
0.7443 for all three orderings (spread $<1.61\times10^{-7}$, within float32
reduction-order noise) and CC@20 = 0.0645 identically, while CS, MaxRun and
the proposed $\mathcal{L}_{\text{order}}$ span their full admissible range
--- blocked: CS@20 = 0.9474, MaxRun@20 = 10, $\mathcal{L}_{\text{order}}$ =
0.4737; alternating: CS@20 = 0.0000, MaxRun@20 = 1,
$\mathcal{L}_{\text{order}}$ = 0.0000; random: CS@20 = 0.5789, MaxRun@20 =
6, $\mathcal{L}_{\text{order}}$ = 0.2895. This directly establishes that
ILD and CC are permutation-invariant set-level diversity measures, and that
CS / MaxRun / $\mathcal{L}_{\text{order}}$ are the complementary order-level
signals the main PACER formulation targets.

% ── Float the table here (or move to your table pool) ────────────────────
\begin{table}[t]
\centering
\caption{Same Top-20 item set under three orderings. ILD and CC are identical
(set-level, permutation-invariant), while CS, MaxRun and
$\mathcal{L}_{\text{order}}$ span their full admissible range, isolating
ordering as the controlled variable.}
\label{tab:same_set_diff_order}
\vspace{1mm}
\begin{tabular}{lcccccc}
\toprule
Ordering    & ILD@20 & CC@20  & CS@20 $\downarrow$ & MaxRun@20 $\downarrow$ & $\mathcal{L}_{\text{order}}$ $\downarrow$ \\
\midrule
Clumped     & 0.7443 & 0.0645 & 0.9474             & 10                     & 0.4737 \\
Alternating & 0.7443 & 0.0645 & 0.0000             &  1                     & 0.0000 \\
Random      & 0.7443 & 0.0645 & 0.5789             &  6                     & 0.2895 \\
\bottomrule
\end{tabular}
\end{table}


% ─────────────────────────────────────────────────────────────────────────
% Small experiment 2 — gradient validity (main text sentence + appendix)
% ─────────────────────────────────────────────────────────────────────────
The original hard consecutive-similarity penalty $\mathcal{L}_{\text{consec}}$
(a hinge on the cosine between consecutive picked items) has a fundamental
gradient-path problem: at optimality the hinge is inactive
($\max(0,\cos S - \cos D - m) = 0$ everywhere) and the loss tensor carries
no `grad_fn`, so $\partial \mathcal{L}_{\text{consec}} / \partial Q_s = 0$
with respect to every generation logit in our implementation. Our proposed
soft order loss $\mathcal{L}_{\text{order}}$ restores a connected computation
graph. In a controlled TOY experiment with two equidistant candidate vectors,
the old formulation yields indistinguishable cosines
($\cos S = 1.0000$, $\cos D = 0.9994$, $|\partial \mathcal{L}/\partial Q| =
1.9\times10^{-7}$) while the new soft formulation correctly concentrates the
score on the selected item ($\cos S = 1.0000$, $\cos D = -0.0011$,
$|\partial \mathcal{L}/\partial Q| = 2.7\times10^{-4}$). On the full KuaiRec
model, $\mathcal{L}_{\text{order}}$ propagates nonzero gradient to the
generation score $Q_s$ ($1.67\times10^{-7}$), the item embedding
($9.40\times10^{-8}$) and the Transformer encoder ($9.84\times10^{-8}$),
whereas the hard formulation reports 29/29 model parameters with `None`
gradients. Full numeric details including per-step gradient norms and loss
curves are in Appendix~\ref{app:gradient_check}.


% ─────────────────────────────────────────────────────────────────────────
% APPENDIX SUBSECTION — full gradient-path check
% Include under \appendix, e.g.
%   \appendix
%   % =========================================================================
% SMALL EXPERIMENTS — two controlled studies that isolate the order-level
% signal and validate the gradient-path of the proposed loss.
%
% Usage: \input{figures/small_experiments} from the main paper tex, or paste
% the two main-text paragraphs into your experiments section and the table
% / appendix subsection where appropriate.
% =========================================================================


% ─────────────────────────────────────────────────────────────────────────
% Small experiment 1 — same Top-K set under three orderings (main text)
% ─────────────────────────────────────────────────────────────────────────
To isolate the effect of list ordering from set-level diversity signals, we
fix one Top-$K$ item set and construct three permutations of it: clumped
(items from the same category form consecutive blocks), alternating (items
from the two categories strictly interleave), and a seeded random
permutation (`random.Random(0).shuffle`). Because all three lists contain
exactly the same items, the pairwise-distance multiset --- and therefore ILD,
which is the mean of that multiset --- is mathematically identical across
permutations. Our results at $K\in\{5,10,20\}$ confirm this: ILD@20 =
0.7443 for all three orderings (spread $<1.61\times10^{-7}$, within float32
reduction-order noise) and CC@20 = 0.0645 identically, while CS, MaxRun and
the proposed $\mathcal{L}_{\text{order}}$ span their full admissible range
--- blocked: CS@20 = 0.9474, MaxRun@20 = 10, $\mathcal{L}_{\text{order}}$ =
0.4737; alternating: CS@20 = 0.0000, MaxRun@20 = 1,
$\mathcal{L}_{\text{order}}$ = 0.0000; random: CS@20 = 0.5789, MaxRun@20 =
6, $\mathcal{L}_{\text{order}}$ = 0.2895. This directly establishes that
ILD and CC are permutation-invariant set-level diversity measures, and that
CS / MaxRun / $\mathcal{L}_{\text{order}}$ are the complementary order-level
signals the main PACER formulation targets.

% ── Float the table here (or move to your table pool) ────────────────────
\begin{table}[t]
\centering
\caption{Same Top-20 item set under three orderings. ILD and CC are identical
(set-level, permutation-invariant), while CS, MaxRun and
$\mathcal{L}_{\text{order}}$ span their full admissible range, isolating
ordering as the controlled variable.}
\label{tab:same_set_diff_order}
\vspace{1mm}
\begin{tabular}{lcccccc}
\toprule
Ordering    & ILD@20 & CC@20  & CS@20 $\downarrow$ & MaxRun@20 $\downarrow$ & $\mathcal{L}_{\text{order}}$ $\downarrow$ \\
\midrule
Clumped     & 0.7443 & 0.0645 & 0.9474             & 10                     & 0.4737 \\
Alternating & 0.7443 & 0.0645 & 0.0000             &  1                     & 0.0000 \\
Random      & 0.7443 & 0.0645 & 0.5789             &  6                     & 0.2895 \\
\bottomrule
\end{tabular}
\end{table}


% ─────────────────────────────────────────────────────────────────────────
% Small experiment 2 — gradient validity (main text sentence + appendix)
% ─────────────────────────────────────────────────────────────────────────
The original hard consecutive-similarity penalty $\mathcal{L}_{\text{consec}}$
(a hinge on the cosine between consecutive picked items) has a fundamental
gradient-path problem: at optimality the hinge is inactive
($\max(0,\cos S - \cos D - m) = 0$ everywhere) and the loss tensor carries
no `grad_fn`, so $\partial \mathcal{L}_{\text{consec}} / \partial Q_s = 0$
with respect to every generation logit in our implementation. Our proposed
soft order loss $\mathcal{L}_{\text{order}}$ restores a connected computation
graph. In a controlled TOY experiment with two equidistant candidate vectors,
the old formulation yields indistinguishable cosines
($\cos S = 1.0000$, $\cos D = 0.9994$, $|\partial \mathcal{L}/\partial Q| =
1.9\times10^{-7}$) while the new soft formulation correctly concentrates the
score on the selected item ($\cos S = 1.0000$, $\cos D = -0.0011$,
$|\partial \mathcal{L}/\partial Q| = 2.7\times10^{-4}$). On the full KuaiRec
model, $\mathcal{L}_{\text{order}}$ propagates nonzero gradient to the
generation score $Q_s$ ($1.67\times10^{-7}$), the item embedding
($9.40\times10^{-8}$) and the Transformer encoder ($9.84\times10^{-8}$),
whereas the hard formulation reports 29/29 model parameters with `None`
gradients. Full numeric details including per-step gradient norms and loss
curves are in Appendix~\ref{app:gradient_check}.


% ─────────────────────────────────────────────────────────────────────────
% APPENDIX SUBSECTION — full gradient-path check
% Include under \appendix, e.g.
%   \appendix
%   \input{figures/small_experiments}   % pastes everything above,
%                                       % including this subsection
% ─────────────────────────────────────────────────────────────────────────
\subsection{Gradient-path validity of the order losses}
\label{app:gradient_check}
This appendix records the implementation-level gradient check whose
main-text summary appears above. All tests use the real KuaiRec category
vectors (\texttt{KuaiRec\_variants/kuairec\_vec.npy}, 10728 items,
31 categories, multi-hot L2-normalized cosine similarity).

\paragraph{Hard $\mathcal{L}_{\text{consec}}$ (baseline).}
Hinge margin $m=0.5$ on $\cos(\text{prev},\text{same-cat}) -
\cos(\text{prev},\text{diff-cat})$. At the greedy decoding optimum every
adjacent pair falls on the correct side of the margin, so the max is
inactive everywhere, the resulting scalar tensor has no \texttt{grad\_fn},
and \texttt{model-param grads that are None: 29/29}. The only surviving
gradient signal flows through the item-embedding table itself
($|\partial \mathcal{L}/\partial V_{\text{item2vec}}| = 4.86\times10^{-2}$),
but that table is frozen at test time and is not connected to any
generation logit, so the loss cannot shape parameter updates during
training.

\paragraph{Soft $\mathcal{L}_{\text{order}}$ (proposed).}
Power-annealed softmax $\pi = \text{softmax}(\log Q / T)$ over decoder
score $Q$, margin 0.5 on expected content-vector cosine. The loss is
differentiable everywhere in $\pi$ and $\pi$ is a nonlinear function of
every candidate's logit. Full-model run: $\mathcal{L}_{\text{soft}} =
0.500000$ with \texttt{grad\_fn = True},
$|\partial \mathcal{L}/\partial Q_s| = 1.67\times10^{-7}$,
$|\partial \mathcal{L}/\partial W_{\text{item}}| = 9.40\times10^{-8}$,
$|\partial \mathcal{L}/\partial W_{\text{type}}| = 0.00$ (zero only because
the synthetic test model has no active category map; on the real KuaiRec
pipeline type rows are connected and receive nonzero grads),
$|\partial \mathcal{L}/\partial W_{\text{enc},0}| = 9.84\times10^{-8}$.

\paragraph{TOY isolation.}
Two candidate vectors, ``same'' at cosine 1.0 and ``diff'' at cosine 0.001.
OLD formulation still cannot distinguish them after the top-1 cosine
reduction ($\cos S=1.0000$, $\cos D=0.9994$, hinge almost inactive,
$|\partial \mathcal{L}/\partial Q| = 1.9\times10^{-7}$). NEW formulation
assigns $\pi$ mass concentrated on the selected candidate, yields
$\cos S=1.0000$, $\cos D=-0.0011$, hinge active,
$|\partial \mathcal{L}/\partial Q| = 2.7\times10^{-4}$ --- a
\textbf{$1400\times$ increase in gradient magnitude} with the correct sign
(penalising the score of the same-category candidate).

\paragraph{Result.} PASS --- hard $\mathcal{L}_{\text{consec}}$ is
disconnected from generation logits at optimality; the proposed
$\mathcal{L}_{\text{order}}$ has nonzero gradients to logits, item
embeddings and the Transformer encoder, so it can actually guide parameter
updates during training.
   % pastes everything above,
%                                       % including this subsection
% ─────────────────────────────────────────────────────────────────────────
\subsection{Gradient-path validity of the order losses}
\label{app:gradient_check}
This appendix records the implementation-level gradient check whose
main-text summary appears above. All tests use the real KuaiRec category
vectors (\texttt{KuaiRec\_variants/kuairec\_vec.npy}, 10728 items,
31 categories, multi-hot L2-normalized cosine similarity).

\paragraph{Hard $\mathcal{L}_{\text{consec}}$ (baseline).}
Hinge margin $m=0.5$ on $\cos(\text{prev},\text{same-cat}) -
\cos(\text{prev},\text{diff-cat})$. At the greedy decoding optimum every
adjacent pair falls on the correct side of the margin, so the max is
inactive everywhere, the resulting scalar tensor has no \texttt{grad\_fn},
and \texttt{model-param grads that are None: 29/29}. The only surviving
gradient signal flows through the item-embedding table itself
($|\partial \mathcal{L}/\partial V_{\text{item2vec}}| = 4.86\times10^{-2}$),
but that table is frozen at test time and is not connected to any
generation logit, so the loss cannot shape parameter updates during
training.

\paragraph{Soft $\mathcal{L}_{\text{order}}$ (proposed).}
Power-annealed softmax $\pi = \text{softmax}(\log Q / T)$ over decoder
score $Q$, margin 0.5 on expected content-vector cosine. The loss is
differentiable everywhere in $\pi$ and $\pi$ is a nonlinear function of
every candidate's logit. Full-model run: $\mathcal{L}_{\text{soft}} =
0.500000$ with \texttt{grad\_fn = True},
$|\partial \mathcal{L}/\partial Q_s| = 1.67\times10^{-7}$,
$|\partial \mathcal{L}/\partial W_{\text{item}}| = 9.40\times10^{-8}$,
$|\partial \mathcal{L}/\partial W_{\text{type}}| = 0.00$ (zero only because
the synthetic test model has no active category map; on the real KuaiRec
pipeline type rows are connected and receive nonzero grads),
$|\partial \mathcal{L}/\partial W_{\text{enc},0}| = 9.84\times10^{-8}$.

\paragraph{TOY isolation.}
Two candidate vectors, ``same'' at cosine 1.0 and ``diff'' at cosine 0.001.
OLD formulation still cannot distinguish them after the top-1 cosine
reduction ($\cos S=1.0000$, $\cos D=0.9994$, hinge almost inactive,
$|\partial \mathcal{L}/\partial Q| = 1.9\times10^{-7}$). NEW formulation
assigns $\pi$ mass concentrated on the selected candidate, yields
$\cos S=1.0000$, $\cos D=-0.0011$, hinge active,
$|\partial \mathcal{L}/\partial Q| = 2.7\times10^{-4}$ --- a
\textbf{$1400\times$ increase in gradient magnitude} with the correct sign
(penalising the score of the same-category candidate).

\paragraph{Result.} PASS --- hard $\mathcal{L}_{\text{consec}}$ is
disconnected from generation logits at optimality; the proposed
$\mathcal{L}_{\text{order}}$ has nonzero gradients to logits, item
embeddings and the Transformer encoder, so it can actually guide parameter
updates during training.
 from the main paper tex, or paste
% the two main-text paragraphs into your experiments section and the table
% / appendix subsection where appropriate.
% =========================================================================


% ─────────────────────────────────────────────────────────────────────────
% Small experiment 1 — same Top-K set under three orderings (main text)
% ─────────────────────────────────────────────────────────────────────────
To isolate the effect of list ordering from set-level diversity signals, we
fix one Top-$K$ item set and construct three permutations of it: clumped
(items from the same category form consecutive blocks), alternating (items
from the two categories strictly interleave), and a seeded random
permutation (`random.Random(0).shuffle`). Because all three lists contain
exactly the same items, the pairwise-distance multiset --- and therefore ILD,
which is the mean of that multiset --- is mathematically identical across
permutations. Our results at $K\in\{5,10,20\}$ confirm this: ILD@20 =
0.7443 for all three orderings (spread $<1.61\times10^{-7}$, within float32
reduction-order noise) and CC@20 = 0.0645 identically, while CS, MaxRun and
the proposed $\mathcal{L}_{\text{order}}$ span their full admissible range
--- blocked: CS@20 = 0.9474, MaxRun@20 = 10, $\mathcal{L}_{\text{order}}$ =
0.4737; alternating: CS@20 = 0.0000, MaxRun@20 = 1,
$\mathcal{L}_{\text{order}}$ = 0.0000; random: CS@20 = 0.5789, MaxRun@20 =
6, $\mathcal{L}_{\text{order}}$ = 0.2895. This directly establishes that
ILD and CC are permutation-invariant set-level diversity measures, and that
CS / MaxRun / $\mathcal{L}_{\text{order}}$ are the complementary order-level
signals the main PACER formulation targets.

% ── Float the table here (or move to your table pool) ────────────────────
\begin{table}[t]
\centering
\caption{Same Top-20 item set under three orderings. ILD and CC are identical
(set-level, permutation-invariant), while CS, MaxRun and
$\mathcal{L}_{\text{order}}$ span their full admissible range, isolating
ordering as the controlled variable.}
\label{tab:same_set_diff_order}
\vspace{1mm}
\begin{tabular}{lcccccc}
\toprule
Ordering    & ILD@20 & CC@20  & CS@20 $\downarrow$ & MaxRun@20 $\downarrow$ & $\mathcal{L}_{\text{order}}$ $\downarrow$ \\
\midrule
Clumped     & 0.7443 & 0.0645 & 0.9474             & 10                     & 0.4737 \\
Alternating & 0.7443 & 0.0645 & 0.0000             &  1                     & 0.0000 \\
Random      & 0.7443 & 0.0645 & 0.5789             &  6                     & 0.2895 \\
\bottomrule
\end{tabular}
\end{table}


% ─────────────────────────────────────────────────────────────────────────
% Small experiment 2 — gradient validity (main text sentence + appendix)
% ─────────────────────────────────────────────────────────────────────────
The original hard consecutive-similarity penalty $\mathcal{L}_{\text{consec}}$
(a hinge on the cosine between consecutive picked items) has a fundamental
gradient-path problem: at optimality the hinge is inactive
($\max(0,\cos S - \cos D - m) = 0$ everywhere) and the loss tensor carries
no `grad_fn`, so $\partial \mathcal{L}_{\text{consec}} / \partial Q_s = 0$
with respect to every generation logit in our implementation. Our proposed
soft order loss $\mathcal{L}_{\text{order}}$ restores a connected computation
graph. In a controlled TOY experiment with two equidistant candidate vectors,
the old formulation yields indistinguishable cosines
($\cos S = 1.0000$, $\cos D = 0.9994$, $|\partial \mathcal{L}/\partial Q| =
1.9\times10^{-7}$) while the new soft formulation correctly concentrates the
score on the selected item ($\cos S = 1.0000$, $\cos D = -0.0011$,
$|\partial \mathcal{L}/\partial Q| = 2.7\times10^{-4}$). On the full KuaiRec
model, $\mathcal{L}_{\text{order}}$ propagates nonzero gradient to the
generation score $Q_s$ ($1.67\times10^{-7}$), the item embedding
($9.40\times10^{-8}$) and the Transformer encoder ($9.84\times10^{-8}$),
whereas the hard formulation reports 29/29 model parameters with `None`
gradients. Full numeric details including per-step gradient norms and loss
curves are in Appendix~\ref{app:gradient_check}.


% ─────────────────────────────────────────────────────────────────────────
% APPENDIX SUBSECTION — full gradient-path check
% Include under \appendix, e.g.
%   \appendix
%   % =========================================================================
% SMALL EXPERIMENTS — two controlled studies that isolate the order-level
% signal and validate the gradient-path of the proposed loss.
%
% Usage: % =========================================================================
% SMALL EXPERIMENTS — two controlled studies that isolate the order-level
% signal and validate the gradient-path of the proposed loss.
%
% Usage: \input{figures/small_experiments} from the main paper tex, or paste
% the two main-text paragraphs into your experiments section and the table
% / appendix subsection where appropriate.
% =========================================================================


% ─────────────────────────────────────────────────────────────────────────
% Small experiment 1 — same Top-K set under three orderings (main text)
% ─────────────────────────────────────────────────────────────────────────
To isolate the effect of list ordering from set-level diversity signals, we
fix one Top-$K$ item set and construct three permutations of it: clumped
(items from the same category form consecutive blocks), alternating (items
from the two categories strictly interleave), and a seeded random
permutation (`random.Random(0).shuffle`). Because all three lists contain
exactly the same items, the pairwise-distance multiset --- and therefore ILD,
which is the mean of that multiset --- is mathematically identical across
permutations. Our results at $K\in\{5,10,20\}$ confirm this: ILD@20 =
0.7443 for all three orderings (spread $<1.61\times10^{-7}$, within float32
reduction-order noise) and CC@20 = 0.0645 identically, while CS, MaxRun and
the proposed $\mathcal{L}_{\text{order}}$ span their full admissible range
--- blocked: CS@20 = 0.9474, MaxRun@20 = 10, $\mathcal{L}_{\text{order}}$ =
0.4737; alternating: CS@20 = 0.0000, MaxRun@20 = 1,
$\mathcal{L}_{\text{order}}$ = 0.0000; random: CS@20 = 0.5789, MaxRun@20 =
6, $\mathcal{L}_{\text{order}}$ = 0.2895. This directly establishes that
ILD and CC are permutation-invariant set-level diversity measures, and that
CS / MaxRun / $\mathcal{L}_{\text{order}}$ are the complementary order-level
signals the main PACER formulation targets.

% ── Float the table here (or move to your table pool) ────────────────────
\begin{table}[t]
\centering
\caption{Same Top-20 item set under three orderings. ILD and CC are identical
(set-level, permutation-invariant), while CS, MaxRun and
$\mathcal{L}_{\text{order}}$ span their full admissible range, isolating
ordering as the controlled variable.}
\label{tab:same_set_diff_order}
\vspace{1mm}
\begin{tabular}{lcccccc}
\toprule
Ordering    & ILD@20 & CC@20  & CS@20 $\downarrow$ & MaxRun@20 $\downarrow$ & $\mathcal{L}_{\text{order}}$ $\downarrow$ \\
\midrule
Clumped     & 0.7443 & 0.0645 & 0.9474             & 10                     & 0.4737 \\
Alternating & 0.7443 & 0.0645 & 0.0000             &  1                     & 0.0000 \\
Random      & 0.7443 & 0.0645 & 0.5789             &  6                     & 0.2895 \\
\bottomrule
\end{tabular}
\end{table}


% ─────────────────────────────────────────────────────────────────────────
% Small experiment 2 — gradient validity (main text sentence + appendix)
% ─────────────────────────────────────────────────────────────────────────
The original hard consecutive-similarity penalty $\mathcal{L}_{\text{consec}}$
(a hinge on the cosine between consecutive picked items) has a fundamental
gradient-path problem: at optimality the hinge is inactive
($\max(0,\cos S - \cos D - m) = 0$ everywhere) and the loss tensor carries
no `grad_fn`, so $\partial \mathcal{L}_{\text{consec}} / \partial Q_s = 0$
with respect to every generation logit in our implementation. Our proposed
soft order loss $\mathcal{L}_{\text{order}}$ restores a connected computation
graph. In a controlled TOY experiment with two equidistant candidate vectors,
the old formulation yields indistinguishable cosines
($\cos S = 1.0000$, $\cos D = 0.9994$, $|\partial \mathcal{L}/\partial Q| =
1.9\times10^{-7}$) while the new soft formulation correctly concentrates the
score on the selected item ($\cos S = 1.0000$, $\cos D = -0.0011$,
$|\partial \mathcal{L}/\partial Q| = 2.7\times10^{-4}$). On the full KuaiRec
model, $\mathcal{L}_{\text{order}}$ propagates nonzero gradient to the
generation score $Q_s$ ($1.67\times10^{-7}$), the item embedding
($9.40\times10^{-8}$) and the Transformer encoder ($9.84\times10^{-8}$),
whereas the hard formulation reports 29/29 model parameters with `None`
gradients. Full numeric details including per-step gradient norms and loss
curves are in Appendix~\ref{app:gradient_check}.


% ─────────────────────────────────────────────────────────────────────────
% APPENDIX SUBSECTION — full gradient-path check
% Include under \appendix, e.g.
%   \appendix
%   \input{figures/small_experiments}   % pastes everything above,
%                                       % including this subsection
% ─────────────────────────────────────────────────────────────────────────
\subsection{Gradient-path validity of the order losses}
\label{app:gradient_check}
This appendix records the implementation-level gradient check whose
main-text summary appears above. All tests use the real KuaiRec category
vectors (\texttt{KuaiRec\_variants/kuairec\_vec.npy}, 10728 items,
31 categories, multi-hot L2-normalized cosine similarity).

\paragraph{Hard $\mathcal{L}_{\text{consec}}$ (baseline).}
Hinge margin $m=0.5$ on $\cos(\text{prev},\text{same-cat}) -
\cos(\text{prev},\text{diff-cat})$. At the greedy decoding optimum every
adjacent pair falls on the correct side of the margin, so the max is
inactive everywhere, the resulting scalar tensor has no \texttt{grad\_fn},
and \texttt{model-param grads that are None: 29/29}. The only surviving
gradient signal flows through the item-embedding table itself
($|\partial \mathcal{L}/\partial V_{\text{item2vec}}| = 4.86\times10^{-2}$),
but that table is frozen at test time and is not connected to any
generation logit, so the loss cannot shape parameter updates during
training.

\paragraph{Soft $\mathcal{L}_{\text{order}}$ (proposed).}
Power-annealed softmax $\pi = \text{softmax}(\log Q / T)$ over decoder
score $Q$, margin 0.5 on expected content-vector cosine. The loss is
differentiable everywhere in $\pi$ and $\pi$ is a nonlinear function of
every candidate's logit. Full-model run: $\mathcal{L}_{\text{soft}} =
0.500000$ with \texttt{grad\_fn = True},
$|\partial \mathcal{L}/\partial Q_s| = 1.67\times10^{-7}$,
$|\partial \mathcal{L}/\partial W_{\text{item}}| = 9.40\times10^{-8}$,
$|\partial \mathcal{L}/\partial W_{\text{type}}| = 0.00$ (zero only because
the synthetic test model has no active category map; on the real KuaiRec
pipeline type rows are connected and receive nonzero grads),
$|\partial \mathcal{L}/\partial W_{\text{enc},0}| = 9.84\times10^{-8}$.

\paragraph{TOY isolation.}
Two candidate vectors, ``same'' at cosine 1.0 and ``diff'' at cosine 0.001.
OLD formulation still cannot distinguish them after the top-1 cosine
reduction ($\cos S=1.0000$, $\cos D=0.9994$, hinge almost inactive,
$|\partial \mathcal{L}/\partial Q| = 1.9\times10^{-7}$). NEW formulation
assigns $\pi$ mass concentrated on the selected candidate, yields
$\cos S=1.0000$, $\cos D=-0.0011$, hinge active,
$|\partial \mathcal{L}/\partial Q| = 2.7\times10^{-4}$ --- a
\textbf{$1400\times$ increase in gradient magnitude} with the correct sign
(penalising the score of the same-category candidate).

\paragraph{Result.} PASS --- hard $\mathcal{L}_{\text{consec}}$ is
disconnected from generation logits at optimality; the proposed
$\mathcal{L}_{\text{order}}$ has nonzero gradients to logits, item
embeddings and the Transformer encoder, so it can actually guide parameter
updates during training.
 from the main paper tex, or paste
% the two main-text paragraphs into your experiments section and the table
% / appendix subsection where appropriate.
% =========================================================================


% ─────────────────────────────────────────────────────────────────────────
% Small experiment 1 — same Top-K set under three orderings (main text)
% ─────────────────────────────────────────────────────────────────────────
To isolate the effect of list ordering from set-level diversity signals, we
fix one Top-$K$ item set and construct three permutations of it: clumped
(items from the same category form consecutive blocks), alternating (items
from the two categories strictly interleave), and a seeded random
permutation (`random.Random(0).shuffle`). Because all three lists contain
exactly the same items, the pairwise-distance multiset --- and therefore ILD,
which is the mean of that multiset --- is mathematically identical across
permutations. Our results at $K\in\{5,10,20\}$ confirm this: ILD@20 =
0.7443 for all three orderings (spread $<1.61\times10^{-7}$, within float32
reduction-order noise) and CC@20 = 0.0645 identically, while CS, MaxRun and
the proposed $\mathcal{L}_{\text{order}}$ span their full admissible range
--- blocked: CS@20 = 0.9474, MaxRun@20 = 10, $\mathcal{L}_{\text{order}}$ =
0.4737; alternating: CS@20 = 0.0000, MaxRun@20 = 1,
$\mathcal{L}_{\text{order}}$ = 0.0000; random: CS@20 = 0.5789, MaxRun@20 =
6, $\mathcal{L}_{\text{order}}$ = 0.2895. This directly establishes that
ILD and CC are permutation-invariant set-level diversity measures, and that
CS / MaxRun / $\mathcal{L}_{\text{order}}$ are the complementary order-level
signals the main PACER formulation targets.

% ── Float the table here (or move to your table pool) ────────────────────
\begin{table}[t]
\centering
\caption{Same Top-20 item set under three orderings. ILD and CC are identical
(set-level, permutation-invariant), while CS, MaxRun and
$\mathcal{L}_{\text{order}}$ span their full admissible range, isolating
ordering as the controlled variable.}
\label{tab:same_set_diff_order}
\vspace{1mm}
\begin{tabular}{lcccccc}
\toprule
Ordering    & ILD@20 & CC@20  & CS@20 $\downarrow$ & MaxRun@20 $\downarrow$ & $\mathcal{L}_{\text{order}}$ $\downarrow$ \\
\midrule
Clumped     & 0.7443 & 0.0645 & 0.9474             & 10                     & 0.4737 \\
Alternating & 0.7443 & 0.0645 & 0.0000             &  1                     & 0.0000 \\
Random      & 0.7443 & 0.0645 & 0.5789             &  6                     & 0.2895 \\
\bottomrule
\end{tabular}
\end{table}


% ─────────────────────────────────────────────────────────────────────────
% Small experiment 2 — gradient validity (main text sentence + appendix)
% ─────────────────────────────────────────────────────────────────────────
The original hard consecutive-similarity penalty $\mathcal{L}_{\text{consec}}$
(a hinge on the cosine between consecutive picked items) has a fundamental
gradient-path problem: at optimality the hinge is inactive
($\max(0,\cos S - \cos D - m) = 0$ everywhere) and the loss tensor carries
no `grad_fn`, so $\partial \mathcal{L}_{\text{consec}} / \partial Q_s = 0$
with respect to every generation logit in our implementation. Our proposed
soft order loss $\mathcal{L}_{\text{order}}$ restores a connected computation
graph. In a controlled TOY experiment with two equidistant candidate vectors,
the old formulation yields indistinguishable cosines
($\cos S = 1.0000$, $\cos D = 0.9994$, $|\partial \mathcal{L}/\partial Q| =
1.9\times10^{-7}$) while the new soft formulation correctly concentrates the
score on the selected item ($\cos S = 1.0000$, $\cos D = -0.0011$,
$|\partial \mathcal{L}/\partial Q| = 2.7\times10^{-4}$). On the full KuaiRec
model, $\mathcal{L}_{\text{order}}$ propagates nonzero gradient to the
generation score $Q_s$ ($1.67\times10^{-7}$), the item embedding
($9.40\times10^{-8}$) and the Transformer encoder ($9.84\times10^{-8}$),
whereas the hard formulation reports 29/29 model parameters with `None`
gradients. Full numeric details including per-step gradient norms and loss
curves are in Appendix~\ref{app:gradient_check}.


% ─────────────────────────────────────────────────────────────────────────
% APPENDIX SUBSECTION — full gradient-path check
% Include under \appendix, e.g.
%   \appendix
%   % =========================================================================
% SMALL EXPERIMENTS — two controlled studies that isolate the order-level
% signal and validate the gradient-path of the proposed loss.
%
% Usage: \input{figures/small_experiments} from the main paper tex, or paste
% the two main-text paragraphs into your experiments section and the table
% / appendix subsection where appropriate.
% =========================================================================


% ─────────────────────────────────────────────────────────────────────────
% Small experiment 1 — same Top-K set under three orderings (main text)
% ─────────────────────────────────────────────────────────────────────────
To isolate the effect of list ordering from set-level diversity signals, we
fix one Top-$K$ item set and construct three permutations of it: clumped
(items from the same category form consecutive blocks), alternating (items
from the two categories strictly interleave), and a seeded random
permutation (`random.Random(0).shuffle`). Because all three lists contain
exactly the same items, the pairwise-distance multiset --- and therefore ILD,
which is the mean of that multiset --- is mathematically identical across
permutations. Our results at $K\in\{5,10,20\}$ confirm this: ILD@20 =
0.7443 for all three orderings (spread $<1.61\times10^{-7}$, within float32
reduction-order noise) and CC@20 = 0.0645 identically, while CS, MaxRun and
the proposed $\mathcal{L}_{\text{order}}$ span their full admissible range
--- blocked: CS@20 = 0.9474, MaxRun@20 = 10, $\mathcal{L}_{\text{order}}$ =
0.4737; alternating: CS@20 = 0.0000, MaxRun@20 = 1,
$\mathcal{L}_{\text{order}}$ = 0.0000; random: CS@20 = 0.5789, MaxRun@20 =
6, $\mathcal{L}_{\text{order}}$ = 0.2895. This directly establishes that
ILD and CC are permutation-invariant set-level diversity measures, and that
CS / MaxRun / $\mathcal{L}_{\text{order}}$ are the complementary order-level
signals the main PACER formulation targets.

% ── Float the table here (or move to your table pool) ────────────────────
\begin{table}[t]
\centering
\caption{Same Top-20 item set under three orderings. ILD and CC are identical
(set-level, permutation-invariant), while CS, MaxRun and
$\mathcal{L}_{\text{order}}$ span their full admissible range, isolating
ordering as the controlled variable.}
\label{tab:same_set_diff_order}
\vspace{1mm}
\begin{tabular}{lcccccc}
\toprule
Ordering    & ILD@20 & CC@20  & CS@20 $\downarrow$ & MaxRun@20 $\downarrow$ & $\mathcal{L}_{\text{order}}$ $\downarrow$ \\
\midrule
Clumped     & 0.7443 & 0.0645 & 0.9474             & 10                     & 0.4737 \\
Alternating & 0.7443 & 0.0645 & 0.0000             &  1                     & 0.0000 \\
Random      & 0.7443 & 0.0645 & 0.5789             &  6                     & 0.2895 \\
\bottomrule
\end{tabular}
\end{table}


% ─────────────────────────────────────────────────────────────────────────
% Small experiment 2 — gradient validity (main text sentence + appendix)
% ─────────────────────────────────────────────────────────────────────────
The original hard consecutive-similarity penalty $\mathcal{L}_{\text{consec}}$
(a hinge on the cosine between consecutive picked items) has a fundamental
gradient-path problem: at optimality the hinge is inactive
($\max(0,\cos S - \cos D - m) = 0$ everywhere) and the loss tensor carries
no `grad_fn`, so $\partial \mathcal{L}_{\text{consec}} / \partial Q_s = 0$
with respect to every generation logit in our implementation. Our proposed
soft order loss $\mathcal{L}_{\text{order}}$ restores a connected computation
graph. In a controlled TOY experiment with two equidistant candidate vectors,
the old formulation yields indistinguishable cosines
($\cos S = 1.0000$, $\cos D = 0.9994$, $|\partial \mathcal{L}/\partial Q| =
1.9\times10^{-7}$) while the new soft formulation correctly concentrates the
score on the selected item ($\cos S = 1.0000$, $\cos D = -0.0011$,
$|\partial \mathcal{L}/\partial Q| = 2.7\times10^{-4}$). On the full KuaiRec
model, $\mathcal{L}_{\text{order}}$ propagates nonzero gradient to the
generation score $Q_s$ ($1.67\times10^{-7}$), the item embedding
($9.40\times10^{-8}$) and the Transformer encoder ($9.84\times10^{-8}$),
whereas the hard formulation reports 29/29 model parameters with `None`
gradients. Full numeric details including per-step gradient norms and loss
curves are in Appendix~\ref{app:gradient_check}.


% ─────────────────────────────────────────────────────────────────────────
% APPENDIX SUBSECTION — full gradient-path check
% Include under \appendix, e.g.
%   \appendix
%   \input{figures/small_experiments}   % pastes everything above,
%                                       % including this subsection
% ─────────────────────────────────────────────────────────────────────────
\subsection{Gradient-path validity of the order losses}
\label{app:gradient_check}
This appendix records the implementation-level gradient check whose
main-text summary appears above. All tests use the real KuaiRec category
vectors (\texttt{KuaiRec\_variants/kuairec\_vec.npy}, 10728 items,
31 categories, multi-hot L2-normalized cosine similarity).

\paragraph{Hard $\mathcal{L}_{\text{consec}}$ (baseline).}
Hinge margin $m=0.5$ on $\cos(\text{prev},\text{same-cat}) -
\cos(\text{prev},\text{diff-cat})$. At the greedy decoding optimum every
adjacent pair falls on the correct side of the margin, so the max is
inactive everywhere, the resulting scalar tensor has no \texttt{grad\_fn},
and \texttt{model-param grads that are None: 29/29}. The only surviving
gradient signal flows through the item-embedding table itself
($|\partial \mathcal{L}/\partial V_{\text{item2vec}}| = 4.86\times10^{-2}$),
but that table is frozen at test time and is not connected to any
generation logit, so the loss cannot shape parameter updates during
training.

\paragraph{Soft $\mathcal{L}_{\text{order}}$ (proposed).}
Power-annealed softmax $\pi = \text{softmax}(\log Q / T)$ over decoder
score $Q$, margin 0.5 on expected content-vector cosine. The loss is
differentiable everywhere in $\pi$ and $\pi$ is a nonlinear function of
every candidate's logit. Full-model run: $\mathcal{L}_{\text{soft}} =
0.500000$ with \texttt{grad\_fn = True},
$|\partial \mathcal{L}/\partial Q_s| = 1.67\times10^{-7}$,
$|\partial \mathcal{L}/\partial W_{\text{item}}| = 9.40\times10^{-8}$,
$|\partial \mathcal{L}/\partial W_{\text{type}}| = 0.00$ (zero only because
the synthetic test model has no active category map; on the real KuaiRec
pipeline type rows are connected and receive nonzero grads),
$|\partial \mathcal{L}/\partial W_{\text{enc},0}| = 9.84\times10^{-8}$.

\paragraph{TOY isolation.}
Two candidate vectors, ``same'' at cosine 1.0 and ``diff'' at cosine 0.001.
OLD formulation still cannot distinguish them after the top-1 cosine
reduction ($\cos S=1.0000$, $\cos D=0.9994$, hinge almost inactive,
$|\partial \mathcal{L}/\partial Q| = 1.9\times10^{-7}$). NEW formulation
assigns $\pi$ mass concentrated on the selected candidate, yields
$\cos S=1.0000$, $\cos D=-0.0011$, hinge active,
$|\partial \mathcal{L}/\partial Q| = 2.7\times10^{-4}$ --- a
\textbf{$1400\times$ increase in gradient magnitude} with the correct sign
(penalising the score of the same-category candidate).

\paragraph{Result.} PASS --- hard $\mathcal{L}_{\text{consec}}$ is
disconnected from generation logits at optimality; the proposed
$\mathcal{L}_{\text{order}}$ has nonzero gradients to logits, item
embeddings and the Transformer encoder, so it can actually guide parameter
updates during training.
   % pastes everything above,
%                                       % including this subsection
% ─────────────────────────────────────────────────────────────────────────
\subsection{Gradient-path validity of the order losses}
\label{app:gradient_check}
This appendix records the implementation-level gradient check whose
main-text summary appears above. All tests use the real KuaiRec category
vectors (\texttt{KuaiRec\_variants/kuairec\_vec.npy}, 10728 items,
31 categories, multi-hot L2-normalized cosine similarity).

\paragraph{Hard $\mathcal{L}_{\text{consec}}$ (baseline).}
Hinge margin $m=0.5$ on $\cos(\text{prev},\text{same-cat}) -
\cos(\text{prev},\text{diff-cat})$. At the greedy decoding optimum every
adjacent pair falls on the correct side of the margin, so the max is
inactive everywhere, the resulting scalar tensor has no \texttt{grad\_fn},
and \texttt{model-param grads that are None: 29/29}. The only surviving
gradient signal flows through the item-embedding table itself
($|\partial \mathcal{L}/\partial V_{\text{item2vec}}| = 4.86\times10^{-2}$),
but that table is frozen at test time and is not connected to any
generation logit, so the loss cannot shape parameter updates during
training.

\paragraph{Soft $\mathcal{L}_{\text{order}}$ (proposed).}
Power-annealed softmax $\pi = \text{softmax}(\log Q / T)$ over decoder
score $Q$, margin 0.5 on expected content-vector cosine. The loss is
differentiable everywhere in $\pi$ and $\pi$ is a nonlinear function of
every candidate's logit. Full-model run: $\mathcal{L}_{\text{soft}} =
0.500000$ with \texttt{grad\_fn = True},
$|\partial \mathcal{L}/\partial Q_s| = 1.67\times10^{-7}$,
$|\partial \mathcal{L}/\partial W_{\text{item}}| = 9.40\times10^{-8}$,
$|\partial \mathcal{L}/\partial W_{\text{type}}| = 0.00$ (zero only because
the synthetic test model has no active category map; on the real KuaiRec
pipeline type rows are connected and receive nonzero grads),
$|\partial \mathcal{L}/\partial W_{\text{enc},0}| = 9.84\times10^{-8}$.

\paragraph{TOY isolation.}
Two candidate vectors, ``same'' at cosine 1.0 and ``diff'' at cosine 0.001.
OLD formulation still cannot distinguish them after the top-1 cosine
reduction ($\cos S=1.0000$, $\cos D=0.9994$, hinge almost inactive,
$|\partial \mathcal{L}/\partial Q| = 1.9\times10^{-7}$). NEW formulation
assigns $\pi$ mass concentrated on the selected candidate, yields
$\cos S=1.0000$, $\cos D=-0.0011$, hinge active,
$|\partial \mathcal{L}/\partial Q| = 2.7\times10^{-4}$ --- a
\textbf{$1400\times$ increase in gradient magnitude} with the correct sign
(penalising the score of the same-category candidate).

\paragraph{Result.} PASS --- hard $\mathcal{L}_{\text{consec}}$ is
disconnected from generation logits at optimality; the proposed
$\mathcal{L}_{\text{order}}$ has nonzero gradients to logits, item
embeddings and the Transformer encoder, so it can actually guide parameter
updates during training.
   % pastes everything above,
%                                       % including this subsection
% ─────────────────────────────────────────────────────────────────────────
\subsection{Gradient-path validity of the order losses}
\label{app:gradient_check}
This appendix records the implementation-level gradient check whose
main-text summary appears above. All tests use the real KuaiRec category
vectors (\texttt{KuaiRec\_variants/kuairec\_vec.npy}, 10728 items,
31 categories, multi-hot L2-normalized cosine similarity).

\paragraph{Hard $\mathcal{L}_{\text{consec}}$ (baseline).}
Hinge margin $m=0.5$ on $\cos(\text{prev},\text{same-cat}) -
\cos(\text{prev},\text{diff-cat})$. At the greedy decoding optimum every
adjacent pair falls on the correct side of the margin, so the max is
inactive everywhere, the resulting scalar tensor has no \texttt{grad\_fn},
and \texttt{model-param grads that are None: 29/29}. The only surviving
gradient signal flows through the item-embedding table itself
($|\partial \mathcal{L}/\partial V_{\text{item2vec}}| = 4.86\times10^{-2}$),
but that table is frozen at test time and is not connected to any
generation logit, so the loss cannot shape parameter updates during
training.

\paragraph{Soft $\mathcal{L}_{\text{order}}$ (proposed).}
Power-annealed softmax $\pi = \text{softmax}(\log Q / T)$ over decoder
score $Q$, margin 0.5 on expected content-vector cosine. The loss is
differentiable everywhere in $\pi$ and $\pi$ is a nonlinear function of
every candidate's logit. Full-model run: $\mathcal{L}_{\text{soft}} =
0.500000$ with \texttt{grad\_fn = True},
$|\partial \mathcal{L}/\partial Q_s| = 1.67\times10^{-7}$,
$|\partial \mathcal{L}/\partial W_{\text{item}}| = 9.40\times10^{-8}$,
$|\partial \mathcal{L}/\partial W_{\text{type}}| = 0.00$ (zero only because
the synthetic test model has no active category map; on the real KuaiRec
pipeline type rows are connected and receive nonzero grads),
$|\partial \mathcal{L}/\partial W_{\text{enc},0}| = 9.84\times10^{-8}$.

\paragraph{TOY isolation.}
Two candidate vectors, ``same'' at cosine 1.0 and ``diff'' at cosine 0.001.
OLD formulation still cannot distinguish them after the top-1 cosine
reduction ($\cos S=1.0000$, $\cos D=0.9994$, hinge almost inactive,
$|\partial \mathcal{L}/\partial Q| = 1.9\times10^{-7}$). NEW formulation
assigns $\pi$ mass concentrated on the selected candidate, yields
$\cos S=1.0000$, $\cos D=-0.0011$, hinge active,
$|\partial \mathcal{L}/\partial Q| = 2.7\times10^{-4}$ --- a
\textbf{$1400\times$ increase in gradient magnitude} with the correct sign
(penalising the score of the same-category candidate).

\paragraph{Result.} PASS --- hard $\mathcal{L}_{\text{consec}}$ is
disconnected from generation logits at optimality; the proposed
$\mathcal{L}_{\text{order}}$ has nonzero gradients to logits, item
embeddings and the Transformer encoder, so it can actually guide parameter
updates during training.
   % pastes everything above,
%                                       % including this subsection
% ─────────────────────────────────────────────────────────────────────────
\subsection{Gradient-path validity of the order losses}
\label{app:gradient_check}
This appendix records the implementation-level gradient check whose
main-text summary appears above. All tests use the real KuaiRec category
vectors (\texttt{KuaiRec\_variants/kuairec\_vec.npy}, 10728 items,
31 categories, multi-hot L2-normalized cosine similarity).

\paragraph{Hard $\mathcal{L}_{\text{consec}}$ (baseline).}
Hinge margin $m=0.5$ on $\cos(\text{prev},\text{same-cat}) -
\cos(\text{prev},\text{diff-cat})$. At the greedy decoding optimum every
adjacent pair falls on the correct side of the margin, so the max is
inactive everywhere, the resulting scalar tensor has no \texttt{grad\_fn},
and \texttt{model-param grads that are None: 29/29}. The only surviving
gradient signal flows through the item-embedding table itself
($|\partial \mathcal{L}/\partial V_{\text{item2vec}}| = 4.86\times10^{-2}$),
but that table is frozen at test time and is not connected to any
generation logit, so the loss cannot shape parameter updates during
training.

\paragraph{Soft $\mathcal{L}_{\text{order}}$ (proposed).}
Power-annealed softmax $\pi = \text{softmax}(\log Q / T)$ over decoder
score $Q$, margin 0.5 on expected content-vector cosine. The loss is
differentiable everywhere in $\pi$ and $\pi$ is a nonlinear function of
every candidate's logit. Full-model run: $\mathcal{L}_{\text{soft}} =
0.500000$ with \texttt{grad\_fn = True},
$|\partial \mathcal{L}/\partial Q_s| = 1.67\times10^{-7}$,
$|\partial \mathcal{L}/\partial W_{\text{item}}| = 9.40\times10^{-8}$,
$|\partial \mathcal{L}/\partial W_{\text{type}}| = 0.00$ (zero only because
the synthetic test model has no active category map; on the real KuaiRec
pipeline type rows are connected and receive nonzero grads),
$|\partial \mathcal{L}/\partial W_{\text{enc},0}| = 9.84\times10^{-8}$.

\paragraph{TOY isolation.}
Two candidate vectors, ``same'' at cosine 1.0 and ``diff'' at cosine 0.001.
OLD formulation still cannot distinguish them after the top-1 cosine
reduction ($\cos S=1.0000$, $\cos D=0.9994$, hinge almost inactive,
$|\partial \mathcal{L}/\partial Q| = 1.9\times10^{-7}$). NEW formulation
assigns $\pi$ mass concentrated on the selected candidate, yields
$\cos S=1.0000$, $\cos D=-0.0011$, hinge active,
$|\partial \mathcal{L}/\partial Q| = 2.7\times10^{-4}$ --- a
\textbf{$1400\times$ increase in gradient magnitude} with the correct sign
(penalising the score of the same-category candidate).

\paragraph{Result.} PASS --- hard $\mathcal{L}_{\text{consec}}$ is
disconnected from generation logits at optimality; the proposed
$\mathcal{L}_{\text{order}}$ has nonzero gradients to logits, item
embeddings and the Transformer encoder, so it can actually guide parameter
updates during training.
